% Options for packages loaded elsewhere
% Options for packages loaded elsewhere
\PassOptionsToPackage{unicode}{hyperref}
\PassOptionsToPackage{hyphens}{url}
\PassOptionsToPackage{dvipsnames,svgnames,x11names}{xcolor}
%
\documentclass[
  letterpaper,
]{report}
\usepackage{xcolor}
\usepackage[lmargin=1in,rmargin=1in]{geometry}
\usepackage{amsmath,amssymb}
\setcounter{secnumdepth}{5}
\usepackage{iftex}
\ifPDFTeX
  \usepackage[T1]{fontenc}
  \usepackage[utf8]{inputenc}
  \usepackage{textcomp} % provide euro and other symbols
\else % if luatex or xetex
  \usepackage{unicode-math} % this also loads fontspec
  \defaultfontfeatures{Scale=MatchLowercase}
  \defaultfontfeatures[\rmfamily]{Ligatures=TeX,Scale=1}
\fi
\usepackage{lmodern}
\ifPDFTeX\else
  % xetex/luatex font selection
\fi
% Use upquote if available, for straight quotes in verbatim environments
\IfFileExists{upquote.sty}{\usepackage{upquote}}{}
\IfFileExists{microtype.sty}{% use microtype if available
  \usepackage[]{microtype}
  \UseMicrotypeSet[protrusion]{basicmath} % disable protrusion for tt fonts
}{}
\makeatletter
\@ifundefined{KOMAClassName}{% if non-KOMA class
  \IfFileExists{parskip.sty}{%
    \usepackage{parskip}
  }{% else
    \setlength{\parindent}{0pt}
    \setlength{\parskip}{6pt plus 2pt minus 1pt}}
}{% if KOMA class
  \KOMAoptions{parskip=half}}
\makeatother
% Make \paragraph and \subparagraph free-standing
\makeatletter
\ifx\paragraph\undefined\else
  \let\oldparagraph\paragraph
  \renewcommand{\paragraph}{
    \@ifstar
      \xxxParagraphStar
      \xxxParagraphNoStar
  }
  \newcommand{\xxxParagraphStar}[1]{\oldparagraph*{#1}\mbox{}}
  \newcommand{\xxxParagraphNoStar}[1]{\oldparagraph{#1}\mbox{}}
\fi
\ifx\subparagraph\undefined\else
  \let\oldsubparagraph\subparagraph
  \renewcommand{\subparagraph}{
    \@ifstar
      \xxxSubParagraphStar
      \xxxSubParagraphNoStar
  }
  \newcommand{\xxxSubParagraphStar}[1]{\oldsubparagraph*{#1}\mbox{}}
  \newcommand{\xxxSubParagraphNoStar}[1]{\oldsubparagraph{#1}\mbox{}}
\fi
\makeatother

\usepackage{color}
\usepackage{fancyvrb}
\newcommand{\VerbBar}{|}
\newcommand{\VERB}{\Verb[commandchars=\\\{\}]}
\DefineVerbatimEnvironment{Highlighting}{Verbatim}{commandchars=\\\{\}}
% Add ',fontsize=\small' for more characters per line
\usepackage{framed}
\definecolor{shadecolor}{RGB}{241,243,245}
\newenvironment{Shaded}{\begin{snugshade}}{\end{snugshade}}
\newcommand{\AlertTok}[1]{\textcolor[rgb]{0.68,0.00,0.00}{#1}}
\newcommand{\AnnotationTok}[1]{\textcolor[rgb]{0.37,0.37,0.37}{#1}}
\newcommand{\AttributeTok}[1]{\textcolor[rgb]{0.40,0.45,0.13}{#1}}
\newcommand{\BaseNTok}[1]{\textcolor[rgb]{0.68,0.00,0.00}{#1}}
\newcommand{\BuiltInTok}[1]{\textcolor[rgb]{0.00,0.23,0.31}{#1}}
\newcommand{\CharTok}[1]{\textcolor[rgb]{0.13,0.47,0.30}{#1}}
\newcommand{\CommentTok}[1]{\textcolor[rgb]{0.37,0.37,0.37}{#1}}
\newcommand{\CommentVarTok}[1]{\textcolor[rgb]{0.37,0.37,0.37}{\textit{#1}}}
\newcommand{\ConstantTok}[1]{\textcolor[rgb]{0.56,0.35,0.01}{#1}}
\newcommand{\ControlFlowTok}[1]{\textcolor[rgb]{0.00,0.23,0.31}{\textbf{#1}}}
\newcommand{\DataTypeTok}[1]{\textcolor[rgb]{0.68,0.00,0.00}{#1}}
\newcommand{\DecValTok}[1]{\textcolor[rgb]{0.68,0.00,0.00}{#1}}
\newcommand{\DocumentationTok}[1]{\textcolor[rgb]{0.37,0.37,0.37}{\textit{#1}}}
\newcommand{\ErrorTok}[1]{\textcolor[rgb]{0.68,0.00,0.00}{#1}}
\newcommand{\ExtensionTok}[1]{\textcolor[rgb]{0.00,0.23,0.31}{#1}}
\newcommand{\FloatTok}[1]{\textcolor[rgb]{0.68,0.00,0.00}{#1}}
\newcommand{\FunctionTok}[1]{\textcolor[rgb]{0.28,0.35,0.67}{#1}}
\newcommand{\ImportTok}[1]{\textcolor[rgb]{0.00,0.46,0.62}{#1}}
\newcommand{\InformationTok}[1]{\textcolor[rgb]{0.37,0.37,0.37}{#1}}
\newcommand{\KeywordTok}[1]{\textcolor[rgb]{0.00,0.23,0.31}{\textbf{#1}}}
\newcommand{\NormalTok}[1]{\textcolor[rgb]{0.00,0.23,0.31}{#1}}
\newcommand{\OperatorTok}[1]{\textcolor[rgb]{0.37,0.37,0.37}{#1}}
\newcommand{\OtherTok}[1]{\textcolor[rgb]{0.00,0.23,0.31}{#1}}
\newcommand{\PreprocessorTok}[1]{\textcolor[rgb]{0.68,0.00,0.00}{#1}}
\newcommand{\RegionMarkerTok}[1]{\textcolor[rgb]{0.00,0.23,0.31}{#1}}
\newcommand{\SpecialCharTok}[1]{\textcolor[rgb]{0.37,0.37,0.37}{#1}}
\newcommand{\SpecialStringTok}[1]{\textcolor[rgb]{0.13,0.47,0.30}{#1}}
\newcommand{\StringTok}[1]{\textcolor[rgb]{0.13,0.47,0.30}{#1}}
\newcommand{\VariableTok}[1]{\textcolor[rgb]{0.07,0.07,0.07}{#1}}
\newcommand{\VerbatimStringTok}[1]{\textcolor[rgb]{0.13,0.47,0.30}{#1}}
\newcommand{\WarningTok}[1]{\textcolor[rgb]{0.37,0.37,0.37}{\textit{#1}}}

\usepackage{longtable,booktabs,array}
\usepackage{calc} % for calculating minipage widths
% Correct order of tables after \paragraph or \subparagraph
\usepackage{etoolbox}
\makeatletter
\patchcmd\longtable{\par}{\if@noskipsec\mbox{}\fi\par}{}{}
\makeatother
% Allow footnotes in longtable head/foot
\IfFileExists{footnotehyper.sty}{\usepackage{footnotehyper}}{\usepackage{footnote}}
\makesavenoteenv{longtable}
\usepackage{graphicx}
\makeatletter
\newsavebox\pandoc@box
\newcommand*\pandocbounded[1]{% scales image to fit in text height/width
  \sbox\pandoc@box{#1}%
  \Gscale@div\@tempa{\textheight}{\dimexpr\ht\pandoc@box+\dp\pandoc@box\relax}%
  \Gscale@div\@tempb{\linewidth}{\wd\pandoc@box}%
  \ifdim\@tempb\p@<\@tempa\p@\let\@tempa\@tempb\fi% select the smaller of both
  \ifdim\@tempa\p@<\p@\scalebox{\@tempa}{\usebox\pandoc@box}%
  \else\usebox{\pandoc@box}%
  \fi%
}
% Set default figure placement to htbp
\def\fps@figure{htbp}
\makeatother





\setlength{\emergencystretch}{3em} % prevent overfull lines

\providecommand{\tightlist}{%
  \setlength{\itemsep}{0pt}\setlength{\parskip}{0pt}}



 


\makeatletter
\@ifpackageloaded{bookmark}{}{\usepackage{bookmark}}
\makeatother
\makeatletter
\@ifpackageloaded{caption}{}{\usepackage{caption}}
\AtBeginDocument{%
\ifdefined\contentsname
  \renewcommand*\contentsname{Table of contents}
\else
  \newcommand\contentsname{Table of contents}
\fi
\ifdefined\listfigurename
  \renewcommand*\listfigurename{List of Figures}
\else
  \newcommand\listfigurename{List of Figures}
\fi
\ifdefined\listtablename
  \renewcommand*\listtablename{List of Tables}
\else
  \newcommand\listtablename{List of Tables}
\fi
\ifdefined\figurename
  \renewcommand*\figurename{Figure}
\else
  \newcommand\figurename{Figure}
\fi
\ifdefined\tablename
  \renewcommand*\tablename{Table}
\else
  \newcommand\tablename{Table}
\fi
}
\@ifpackageloaded{float}{}{\usepackage{float}}
\floatstyle{ruled}
\@ifundefined{c@chapter}{\newfloat{codelisting}{h}{lop}}{\newfloat{codelisting}{h}{lop}[chapter]}
\floatname{codelisting}{Listing}
\newcommand*\listoflistings{\listof{codelisting}{List of Listings}}
\makeatother
\makeatletter
\makeatother
\makeatletter
\@ifpackageloaded{caption}{}{\usepackage{caption}}
\@ifpackageloaded{subcaption}{}{\usepackage{subcaption}}
\makeatother
\usepackage{bookmark}
\IfFileExists{xurl.sty}{\usepackage{xurl}}{} % add URL line breaks if available
\urlstyle{same}
\hypersetup{
  pdftitle={Survey Question Harmonization Constraints},
  pdfauthor={Brock Webb},
  colorlinks=true,
  linkcolor={blue},
  filecolor={Maroon},
  citecolor={Blue},
  urlcolor={Blue},
  pdfcreator={LaTeX via pandoc}}


\title{Survey Question Harmonization Constraints}
\usepackage{etoolbox}
\makeatletter
\providecommand{\subtitle}[1]{% add subtitle to \maketitle
  \apptocmd{\@title}{\par {\large #1 \par}}{}{}
}
\makeatother
\subtitle{AI-Assisted Analysis of Federal Survey Consolidation
Potential}
\author{Brock Webb}
\date{2026-02-02}
\begin{document}
\maketitle

\renewcommand*\contentsname{Table of contents}
{
\hypersetup{linkcolor=}
\setcounter{tocdepth}{2}
\tableofcontents
}

\bookmarksetup{startatroot}

\chapter{Survey Question Harmonization
Constraints}\label{survey-question-harmonization-constraints}

AI-Assisted Analysis of Federal Survey Consolidation Potential

\hfill\break

\bookmarksetup{startatroot}

\chapter*{Preface}\label{preface}
\addcontentsline{toc}{chapter}{Preface}

\markboth{Preface}{Preface}

This report presents findings from an AI-assisted analysis of federal
survey question consolidation potential. Using a multi-model ensemble
approach, we analyzed 1,598 question pairs from the Current Population
Survey (CPS) and FoodAPS against the American Community Survey (ACS) to
determine:

\begin{enumerate}
\def\labelenumi{\arabic{enumi}.}
\tightlist
\item
  \textbf{What proportion of survey questions can be consolidated?}
\item
  \textbf{What barriers prevent consolidation?}
\item
  \textbf{Can AI-assisted methods accelerate this traditionally
  labor-intensive analysis?}
\end{enumerate}

\section{Key Findings}\label{key-findings}

\begin{itemize}
\tightlist
\item
  \textbf{\textasciitilde44\% consolidation potential} identified across
  380 questions
\item
  \textbf{97\% of failures} stem from construct/concept differences (not
  fixable)
\item
  \textbf{AI handles 76\%} of classifications with high confidence,
  routing only 24\% to expert review
\item
  \textbf{Reproducible methodology} applicable to other survey
  harmonization efforts
\end{itemize}

\section{Report Structure}\label{report-structure}

\begin{itemize}
\tightlist
\item
  \textbf{Chapters 1-3}: Background, research context, and harmonization
  framework
\item
  \textbf{Chapters 4-6}: Methodology, results, and discussion
\item
  \textbf{Appendices}: Technical details, taxonomy definitions, and
  expert review tables
\end{itemize}

\section{Data Availability}\label{data-availability}

All analysis code, intermediate outputs, and expert review tables are
available in the project repository:
\texttt{reports/03\_harmonization\_constraints/}

\begin{center}\rule{0.5\linewidth}{0.5pt}\end{center}

\textbf{Note:} This is Report 03 of a multi-phase project analyzing
federal survey questions. Reports 01 (concept classification) and 02
(pairwise comparison setup) established the foundation for this
harmonization analysis.

\bookmarksetup{startatroot}

\chapter*{Executive Summary}\label{executive-summary}
\addcontentsline{toc}{chapter}{Executive Summary}

\markboth{Executive Summary}{Executive Summary}

\section*{Problem}\label{problem}
\addcontentsline{toc}{section}{Problem}

\markright{Problem}

Federal surveys ask overlapping questions across agencies, creating
respondent burden and data integration challenges. Determining which
questions can be harmonized requires expert judgment and traditionally
takes weeks or months of manual review.

\section*{Approach}\label{approach}
\addcontentsline{toc}{section}{Approach}

\markright{Approach}

We developed an AI-assisted pairwise analysis system using: -
\textbf{Multi-model ensemble} (OpenAI, Anthropic, Google) for initial
classification - \textbf{Arbitration pattern} to resolve disagreements -
\textbf{Question-level rollup} using two-axis triage (Borda direction ×
Entropy stability) - \textbf{1,598 question pairs} analyzed (CPS-ACS,
FoodAPS-ACS)

\section*{Key Results}\label{key-results}
\addcontentsline{toc}{section}{Key Results}

\markright{Key Results}

\begin{itemize}
\tightlist
\item
  \textbf{CPS}: 41.7\% consolidable (100 of 240 questions)
\item
  \textbf{FoodAPS}: 48.6\% consolidable (68 of 140 questions)
\item
  \textbf{Combined}: \textasciitilde44\% consolidation potential
\end{itemize}

\textbf{Barrier Analysis:} - 97\% of failures = Construct/Concept (CC)
differences - CC.1 (concept definition) accounts for 70\% of all
incompatible pairs - Not fixable without re-fielding questions

\textbf{Expert Review Load:} - 76\% auto-processed with high confidence
- 24\% (93 questions) routed to expert review - AI handles grunt work;
humans judge edge cases

\section*{Deliverables}\label{deliverables}
\addcontentsline{toc}{section}{Deliverables}

\markright{Deliverables}

\begin{enumerate}
\def\labelenumi{\arabic{enumi}.}
\tightlist
\item
  \textbf{168 consolidable mappings} with specific question-to-question
  recommendations
\item
  \textbf{212 non-consolidable questions} with barrier codes explaining
  why
\item
  \textbf{93 questions flagged for expert review} (uncertain or
  contested cases)
\item
  \textbf{Reproducible methodology} applicable to additional federal
  surveys
\end{enumerate}

\section*{Implications}\label{implications}
\addcontentsline{toc}{section}{Implications}

\markright{Implications}

This work demonstrates that AI-assisted methods can: -
\textbf{Accelerate} survey harmonization analysis (days instead of
weeks) - \textbf{Preserve} human oversight and expert judgment -
\textbf{Focus} expert effort on highest-ROI decisions - \textbf{Scale}
to analyze entire federal survey ecosystem

\section*{Next Steps}\label{next-steps}
\addcontentsline{toc}{section}{Next Steps}

\markright{Next Steps}

\begin{enumerate}
\def\labelenumi{\arabic{enumi}.}
\tightlist
\item
  Expert validation of recommendations
\item
  Expansion to additional federal surveys
\item
  Application to survey editing and imputation workflows
\end{enumerate}

\bookmarksetup{startatroot}

\chapter{Introduction}\label{introduction}

\section{Problem Statement}\label{problem-statement}

The federal statistical system conducts dozens of surveys annually, many
asking overlapping questions about demographics, economics, health, and
social conditions. This redundancy creates:

\begin{enumerate}
\def\labelenumi{\arabic{enumi}.}
\tightlist
\item
  \textbf{Respondent burden}: Same households answering similar
  questions across multiple surveys
\item
  \textbf{Data integration challenges}: Incompatible question wording
  prevents cross-survey analysis
\item
  \textbf{Resource inefficiency}: Agencies duplicating effort in
  questionnaire design and data collection
\end{enumerate}

Determining which questions can be consolidated or harmonized requires
deep subject-matter expertise and traditionally involves weeks or months
of manual review by survey methodologists.

\section{Research Questions}\label{research-questions}

This report addresses three core questions:

\subsection{RQ1: Consolidation
Potential}\label{rq1-consolidation-potential}

\textbf{What proportion of federal survey questions can be consolidated
with existing ACS questions?}

\begin{itemize}
\tightlist
\item
  Hypothesis: Significant consolidation potential exists but requires
  systematic identification
\item
  Approach: Exhaustive pairwise comparison of CPS/FoodAPS questions
  against ACS
\item
  Outcome: Quantify F1 (direct), F2 (with adjustment), F3 (incompatible)
  rates
\end{itemize}

\subsection{RQ2: Barrier Taxonomy}\label{rq2-barrier-taxonomy}

\textbf{What barriers prevent survey question consolidation?}

\begin{itemize}
\tightlist
\item
  Hypothesis: Most barriers stem from construct differences, not
  operational issues
\item
  Approach: Classify failures using DataSHaPER/Maelstrom harmonization
  framework
\item
  Outcome: Identify primary barrier types and their prevalence
\end{itemize}

\subsection{RQ3: AI-Assisted Methods}\label{rq3-ai-assisted-methods}

\textbf{Can AI-assisted methods accelerate this traditionally
labor-intensive analysis?}

\begin{itemize}
\tightlist
\item
  Hypothesis: LLM ensemble can handle routine classifications, reserving
  expert judgment for edge cases
\item
  Approach: Multi-model rating + arbitration + two-axis triage
\item
  Outcome: Measure agreement rates, expert review load reduction
\end{itemize}

\section{Scope}\label{scope}

\subsection{Surveys Analyzed}\label{surveys-analyzed}

\begin{itemize}
\tightlist
\item
  \textbf{American Community Survey (ACS)}: Target survey with
  comprehensive demographic/economic questions
\item
  \textbf{Current Population Survey (CPS)}: Employment-focused survey
  (240 questions)
\item
  \textbf{FoodAPS}: Food security survey (140 questions)
\end{itemize}

Total: \textbf{380 unique source questions} compared against ACS,
yielding \textbf{1,598 question pairs}.

\subsection{Framework Applied}\label{framework-applied}

We adopt the DataSHaPER/Maelstrom retrospective harmonization framework:
- \textbf{F1}: Direct recode (mechanically transformable) - \textbf{F2}:
Statistical adjustment (requires modeling/bridging) - \textbf{F3}:
Incompatible (fundamental barriers prevent harmonization)

\subsection{What This Report Is NOT}\label{what-this-report-is-not}

\begin{itemize}
\tightlist
\item
  Not a recommendation to eliminate any surveys (policy decision)
\item
  Not a claim that all F1/F2 pairs should be consolidated (requires
  stakeholder input)
\item
  Not a replacement for expert judgment (AI assists, experts decide)
\end{itemize}

\section{Prior Work}\label{prior-work}

This is \textbf{Report 03} of a multi-phase project:

\begin{itemize}
\tightlist
\item
  \textbf{Report 01}: Concept classification of 7,400+ questions using
  LLM ensemble
\item
  \textbf{Report 02}: Pairwise comparison setup and sampling strategy
\item
  \textbf{Report 03} (this report): Harmonization constraints and
  consolidation analysis
\end{itemize}

Each report builds on the previous, with validated outputs feeding
forward.

\section{Report Structure}\label{report-structure-1}

The remainder of this report is organized as:

\begin{itemize}
\tightlist
\item
  \textbf{Chapter 2 (Background)}: Survey harmonization frameworks and
  prior work
\item
  \textbf{Chapter 3 (Methodology)}: 5-stage pipeline from rating to
  deliverables
\item
  \textbf{Chapter 4 (Results)}: Consolidation rates, barrier analysis,
  agreement statistics
\item
  \textbf{Chapter 5 (Discussion)}: Interpretation, limitations,
  implications
\item
  \textbf{Chapter 6 (Conclusion)}: Summary and future directions
\item
  \textbf{Appendices}: Taxonomy definitions, methodology decisions,
  expert review tables
\end{itemize}

\section{Key Contributions}\label{key-contributions}

This work contributes:

\begin{enumerate}
\def\labelenumi{\arabic{enumi}.}
\tightlist
\item
  \textbf{Empirical findings} on federal survey consolidation potential
  (\textasciitilde44\%)
\item
  \textbf{Barrier taxonomy application} showing CC dominance (97\% of
  failures)
\item
  \textbf{Operational triage framework} (two-axis: Borda direction ×
  Entropy stability)
\item
  \textbf{Reproducible methodology} for AI-assisted survey harmonization
\end{enumerate}

\begin{center}\rule{0.5\linewidth}{0.5pt}\end{center}

\textbf{Note}: All data, code, and intermediate outputs are available in
\texttt{reports/03\_harmonization\_constraints/} for replication and
extension.

\bookmarksetup{startatroot}

\chapter{Background}\label{background}

\section{Survey Harmonization
Frameworks}\label{survey-harmonization-frameworks}

\subsection{DataSHaPER/Maelstrom
Framework}\label{datashapermaelstrom-framework}

The DataSHaPER (Data Schema and Harmonization Platform for
Epidemiological Research) framework, developed by Fortier et al.~(2011,
2017), provides a structured approach to retrospective data
harmonization. This framework classifies variable pairs into three
feasibility categories:

\begin{longtable}[]{@{}
  >{\raggedright\arraybackslash}p{(\linewidth - 6\tabcolsep) * \real{0.1250}}
  >{\raggedright\arraybackslash}p{(\linewidth - 6\tabcolsep) * \real{0.2708}}
  >{\raggedright\arraybackslash}p{(\linewidth - 6\tabcolsep) * \real{0.2500}}
  >{\raggedright\arraybackslash}p{(\linewidth - 6\tabcolsep) * \real{0.3542}}@{}}
\toprule\noalign{}
\begin{minipage}[b]{\linewidth}\raggedright
Code
\end{minipage} & \begin{minipage}[b]{\linewidth}\raggedright
Feasibility
\end{minipage} & \begin{minipage}[b]{\linewidth}\raggedright
Definition
\end{minipage} & \begin{minipage}[b]{\linewidth}\raggedright
Action Required
\end{minipage} \\
\midrule\noalign{}
\endhead
\bottomrule\noalign{}
\endlastfoot
\textbf{F1} & \textbf{Direct recode} & Variables are mechanically
transformable through simple recoding, collapsing categories, or unit
conversion & Simple data transformation \\
\textbf{F2} & \textbf{Statistical adjustment} & Variables require
modeling, imputation, or bridging studies to make comparable &
Statistical harmonization methods \\
\textbf{F3} & \textbf{Incompatible} & Variables measure fundamentally
different constructs and cannot be harmonized without re-fielding & No
harmonization possible \\
\end{longtable}

\textbf{Source}: Fortier et al.~(2011) applied this framework across 53
epidemiological studies and found that 38\% of variables could be
harmonized through direct or statistical methods.

\subsection{Barrier Taxonomy}\label{barrier-taxonomy}

When variables are classified as F3 (incompatible), the framework
identifies specific barrier types:

\subsubsection{Construct/Concept (CC)}\label{constructconcept-cc}

\begin{itemize}
\tightlist
\item
  \textbf{CC.1}: Concept definition differences (e.g., ``employment''
  including/excluding unpaid work)
\item
  \textbf{CC.2}: Operationalization differences (different behavioral
  indicators for same concept)
\item
  \textbf{CC.3}: Boundary conditions (different thresholds or cutoffs)
\item
  \textbf{CC.4}: Scope inclusions (different components counted)
\end{itemize}

\subsubsection{Temporal Constraints (TC)}\label{temporal-constraints-tc}

\begin{itemize}
\tightlist
\item
  \textbf{TC.1}: Reference period length (7-day vs 12-month recall)
\item
  \textbf{TC.2}: Temporal framing (point-in-time vs habitual vs
  retrospective)
\item
  \textbf{TC.3}: Calendar alignment (fixed vs rolling reference periods)
\end{itemize}

\subsubsection{Response Scale (RS)}\label{response-scale-rs}

\begin{itemize}
\tightlist
\item
  \textbf{RS.1}: Scale type (binary vs Likert vs continuous)
\item
  \textbf{RS.2}: Category structure (different number/boundaries)
\item
  \textbf{RS.3}: Anchoring/labels (different verbal anchors)
\item
  \textbf{RS.4}: Numeric vs verbal scales
\end{itemize}

\subsubsection{Population/Coverage (PC)}\label{populationcoverage-pc}

\begin{itemize}
\tightlist
\item
  \textbf{PC.1}: Universe definition (target population differs)
\item
  \textbf{PC.2}: Frame exclusions (different sampling exclusions)
\item
  \textbf{PC.3}: Age bounds (different age eligibility)
\item
  \textbf{PC.4}: Geographic scope (different coverage)
\end{itemize}

\subsubsection{Mode/Context (MC)}\label{modecontext-mc}

\begin{itemize}
\tightlist
\item
  \textbf{MC.1}: Interview mode (CATI vs web vs in-person)
\item
  \textbf{MC.2}: Question routing (different skip patterns)
\item
  \textbf{MC.3}: Contextual priming (question order effects)
\item
  \textbf{MC.4}: Proxy response rules
\end{itemize}

\subsubsection{Processing/Metadata (PM)}\label{processingmetadata-pm}

\begin{itemize}
\tightlist
\item
  \textbf{PM.1}: Coding schemes (different classification systems)
\item
  \textbf{PM.2}: Derived variables (different construction algorithms)
\item
  \textbf{PM.3}: Documentation gaps (insufficient metadata)
\end{itemize}

\textbf{See Appendix A} for complete taxonomy definitions with examples.

\section{Related Literature}\label{related-literature}

\subsection{Survey Data Harmonization}\label{survey-data-harmonization}

\begin{itemize}
\tightlist
\item
  \textbf{Wolf et al.~(2016)}: Harmonizing survey questions between
  cultures and over time
\item
  \textbf{Saris \& Gallhofer (2014)}: Design, evaluation, and analysis
  of questionnaires for survey research
\item
  \textbf{Slomczynski \& Tomescu-Dubrow (2018)}: Basic principles of
  survey data recycling
\end{itemize}

\subsection{AI-Assisted Survey
Analysis}\label{ai-assisted-survey-analysis}

\section{Prior Work: Reports 01-02}\label{prior-work-reports-01-02}

\subsection{Report 01: Concept
Classification}\label{report-01-concept-classification}

\begin{itemize}
\tightlist
\item
  Mapped 7,400+ survey questions to Census Bureau taxonomy
\item
  Used LLM ensemble (Claude Haiku + GPT-4o-mini) with arbitration
\item
  Achieved 99.5\% categorization success rate
\item
  Cohen's κ = 0.842 (almost perfect agreement)
\end{itemize}

\subsection{Report 02: Pairwise Comparison
Setup}\label{report-02-pairwise-comparison-setup}

\begin{itemize}
\tightlist
\item
  Identified 1,598 question pairs for harmonization analysis
\item
  Selected CPS and FoodAPS as representative source surveys
\item
  Developed sampling strategy based on concept categories
\item
  Established pair-level comparison methodology
\end{itemize}

\subsection{Building Forward}\label{building-forward}

Report 03 (this report) takes the validated pairwise comparisons from
Report 02 and: 1. Classifies each pair using harmonization framework
(F1/F2/F3) 2. Identifies barrier codes for F3 pairs 3. Rolls up
pair-level results to question-level consolidability 4. Generates expert
review tables for stakeholder use

\section{Census Bureau Context}\label{census-bureau-context}

The Census Bureau conducts multiple surveys covering demographics,
economics, housing, and social conditions: - \textbf{ACS}: Comprehensive
demographic/economic data, \textasciitilde3.5M households annually -
\textbf{CPS}: Labor force statistics, \textasciitilde60K households
monthly - \textbf{SIPP}: Household dynamics and program participation -
\textbf{Others}: American Housing Survey, Survey of Income and Program
Participation, etc.

Harmonization across these surveys would enable: - Cross-survey data
integration - Reduced respondent burden through shared questions -
Improved data quality through consistent measurement - Resource
efficiency in questionnaire design

\section{Research Gap}\label{research-gap}

While survey harmonization frameworks exist (DataSHaPER, SDR
methodology), \textbf{no prior work has}: 1. Applied these frameworks
systematically to federal survey questions at scale 2. Used AI-assisted
methods to accelerate the classification process 3. Quantified
consolidation potential across major federal surveys 4. Developed
operational triage frameworks for expert review prioritization

This report addresses these gaps.

\begin{center}\rule{0.5\linewidth}{0.5pt}\end{center}

\textbf{Key Takeaway}: Survey harmonization is well-established in
epidemiology and international survey research, but has not been
systematically applied to the U.S. federal survey ecosystem. This report
demonstrates that AI-assisted methods can scale this analysis.

\bookmarksetup{startatroot}

\chapter{Methodology}\label{methodology}

\section{Overview}\label{overview}

We developed a 5-stage pipeline for AI-assisted survey harmonization
analysis:

\begin{verbatim}
Stage 1: Rating (3-model ensemble)
   ↓
Stage 2: Agreement Analysis (inter-rater reliability)
   ↓
Stage 3: Arbitration (disagreement resolution)
   ↓
Stage 4: Question-Level Rollup (best-match + triage)
   ↓
Stage 5: Deliverables (expert review tables)
\end{verbatim}

All code is available in
\texttt{reports/03\_harmonization\_constraints/scripts/}.

\begin{center}\rule{0.5\linewidth}{0.5pt}\end{center}

\section{Stage 1: Rating}\label{stage-1-rating}

\subsection{Objective}\label{objective}

Classify 1,598 question pairs using harmonization framework (F1/F2/F3 +
barrier codes).

\subsection{Multi-Model Ensemble}\label{multi-model-ensemble}

We employed three frontier LLMs as independent raters:

\begin{longtable}[]{@{}
  >{\raggedright\arraybackslash}p{(\linewidth - 4\tabcolsep) * \real{0.2059}}
  >{\raggedright\arraybackslash}p{(\linewidth - 4\tabcolsep) * \real{0.2647}}
  >{\raggedright\arraybackslash}p{(\linewidth - 4\tabcolsep) * \real{0.5294}}@{}}
\toprule\noalign{}
\begin{minipage}[b]{\linewidth}\raggedright
Model
\end{minipage} & \begin{minipage}[b]{\linewidth}\raggedright
Version
\end{minipage} & \begin{minipage}[b]{\linewidth}\raggedright
Behavior Profile
\end{minipage} \\
\midrule\noalign{}
\endhead
\bottomrule\noalign{}
\endlastfoot
\textbf{OpenAI} & gpt-4o-mini & Moderate synthesis, slight
self-consistency bias \\
\textbf{Anthropic} & claude-3-5-haiku-20241022 & High synthesis,
neutral \\
\textbf{Google} & gemini-2.0-flash-exp & Deferential, conservative \\
\end{longtable}

\textbf{Why three models?} - Reduces single-model bias - Different
training data and architectures - Behavioral diversity improves coverage

\subsection{Prompt Design}\label{prompt-design}

Each model received: 1. \textbf{Harmonization framework definitions}
(F1/F2/F3) 2. \textbf{Barrier taxonomy} (CC, TC, RS, PC, MC, PM) 3.
\textbf{Question pair} (source + target text) 4. \textbf{Task}: Classify
feasibility + provide reasoning

\textbf{Output format} (structured JSON):

\begin{Shaded}
\begin{Highlighting}[]
\FunctionTok{\{}
  \DataTypeTok{"feasibility"}\FunctionTok{:} \StringTok{"F1|F2|F3"}\FunctionTok{,}
  \DataTypeTok{"barrier\_code"}\FunctionTok{:} \StringTok{"CC.1|TC.2|..."} \ErrorTok{(if} \ErrorTok{F3)}\FunctionTok{,}
  \DataTypeTok{"confidence"}\FunctionTok{:} \StringTok{"HIGH|MODERATE|LOW"}\FunctionTok{,}
  \DataTypeTok{"reasoning"}\FunctionTok{:} \StringTok{"Explanation..."}
\FunctionTok{\}}
\end{Highlighting}
\end{Shaded}

\subsection{Execution}\label{execution}

\begin{itemize}
\tightlist
\item
  \textbf{Parallel processing}: 6 workers per model
\item
  \textbf{Batch size}: 10 question pairs per API call
\item
  \textbf{Error handling}: Exponential backoff for rate limits
\item
  \textbf{Checkpointing}: Resume capability for interrupted runs
\item
  \textbf{Output}:
  \texttt{output/results/stage1\_\{model\}\_results.jsonl}
\end{itemize}

\subsection{Data Quality}\label{data-quality}

\begin{longtable}[]{@{}llll@{}}
\toprule\noalign{}
Metric & OpenAI & Anthropic & Google \\
\midrule\noalign{}
\endhead
\bottomrule\noalign{}
\endlastfoot
\textbf{Pairs rated} & 1,598 & 1,598 & 751 (47\%) \\
\textbf{Success rate} & 100\% & 100\% & 47\% \\
\textbf{Avg confidence} & 0.89 & 0.91 & 0.87 \\
\end{longtable}

\textbf{Note}: Google rate-limited at 751 pairs. Remaining analysis uses
OpenAI-Anthropic two-way comparison with Google as validation subset.

\begin{center}\rule{0.5\linewidth}{0.5pt}\end{center}

\section{Stage 2: Agreement Analysis}\label{stage-2-agreement-analysis}

\subsection{Objective}\label{objective-1}

Quantify inter-rater reliability and identify disagreements for
arbitration.

\subsection{Metrics Calculated}\label{metrics-calculated}

\subsubsection{Cohen's Kappa (Pairwise)}\label{cohens-kappa-pairwise}

\begin{itemize}
\tightlist
\item
  \textbf{OpenAI-Anthropic}: κ = 0.845 (almost perfect)
\item
  \textbf{OpenAI-Google}: κ = 0.796 (substantial)
\item
  \textbf{Anthropic-Google}: κ = 0.833 (almost perfect)
\end{itemize}

\subsubsection{Fleiss' Kappa (Three-Way)}\label{fleiss-kappa-three-way}

\begin{itemize}
\tightlist
\item
  \textbf{All three models} (751 pairs with Google): κ = 0.833
\end{itemize}

\textbf{Interpretation}: High agreement validates that harmonization
judgments are consistent across models, even with different
architectures and training data.

\subsection{Agreement Patterns}\label{agreement-patterns}

\begin{longtable}[]{@{}lll@{}}
\toprule\noalign{}
Level & OpenAI-Anthropic & Description \\
\midrule\noalign{}
\endhead
\bottomrule\noalign{}
\endlastfoot
\textbf{Topic (L1)} & 91.2\% & High-level feasibility \\
\textbf{Subtopic (L2)} & 84.5\% & Barrier category \\
\textbf{Full code} & 78.9\% & Complete barrier sub-code \\
\end{longtable}

\textbf{Finding}: Models agree more on broad feasibility than specific
barrier codes, consistent with hierarchical classification difficulty.

\subsection{Disagreement
Identification}\label{disagreement-identification}

\begin{itemize}
\tightlist
\item
  \textbf{Pairs needing arbitration}: 339 (21.2\% of 1,598)
\item
  \textbf{Criteria}: OpenAI ≠ Anthropic on feasibility or barrier code
\item
  \textbf{Output}: \texttt{output/analysis/arbitration\_candidates.csv}
\end{itemize}

\begin{center}\rule{0.5\linewidth}{0.5pt}\end{center}

\section{Stage 3: Arbitration}\label{stage-3-arbitration}

\subsection{Objective}\label{objective-2}

Resolve disagreements through higher-capability model arbitration.

\subsection{Arbitration Models}\label{arbitration-models}

\begin{longtable}[]{@{}
  >{\raggedright\arraybackslash}p{(\linewidth - 4\tabcolsep) * \real{0.3182}}
  >{\raggedright\arraybackslash}p{(\linewidth - 4\tabcolsep) * \real{0.4091}}
  >{\raggedright\arraybackslash}p{(\linewidth - 4\tabcolsep) * \real{0.2727}}@{}}
\toprule\noalign{}
\begin{minipage}[b]{\linewidth}\raggedright
Model
\end{minipage} & \begin{minipage}[b]{\linewidth}\raggedright
Version
\end{minipage} & \begin{minipage}[b]{\linewidth}\raggedright
Role
\end{minipage} \\
\midrule\noalign{}
\endhead
\bottomrule\noalign{}
\endlastfoot
\textbf{OpenAI} & gpt-4o & Arbitrator (stronger than rater) \\
\textbf{Anthropic} & claude-3-5-sonnet-20241022 & Arbitrator \\
\textbf{Google} & gemini-2.0-flash-thinking-exp & Arbitrator
(experimental reasoning) \\
\end{longtable}

\subsection{Arbitration Prompt}\label{arbitration-prompt}

Arbitrators received: 1. \textbf{Source and target question text} 2.
\textbf{Initial ratings} from 2-3 raters (with reasoning) 3.
\textbf{Disagreement summary} 4. \textbf{Task}: Provide final verdict
with detailed justification

\subsection{Voting Logic}\label{voting-logic}

When arbitrators disagree, final verdict determined by: 1.
\textbf{Majority vote} (if 2+ agree) 2. \textbf{Borda count} (ranked
preference if no majority) 3. \textbf{Entropy check} (flag high
disagreement for expert review)

\subsection{Arbitration Results}\label{arbitration-results}

\begin{longtable}[]{@{}ll@{}}
\toprule\noalign{}
Metric & Value \\
\midrule\noalign{}
\endhead
\bottomrule\noalign{}
\endlastfoot
\textbf{Pairs arbitrated} & 339 \\
\textbf{Arbitrator agreement} & 87.6\% \\
\textbf{Confidence distribution} & HIGH: 68\%, MODERATE: 24\%, LOW:
8\% \\
\end{longtable}

\textbf{Quality check} (11 validation tests): ✅ All passed

\subsection{Behavioral Findings}\label{behavioral-findings}

\begin{itemize}
\tightlist
\item
  \textbf{OpenAI}: Tends toward F2 when uncertain (optimistic about
  statistical adjustment)
\item
  \textbf{Anthropic}: Balanced, detailed reasoning
\item
  \textbf{Google}: Conservative, often defers to initial raters
\end{itemize}

\textbf{See}:
\texttt{docs/FINDINGS\_R03\_S3\_001\_arbitration\_analysis.md} for
detailed behavioral analysis.

\begin{center}\rule{0.5\linewidth}{0.5pt}\end{center}

\section{Stage 4: Question-Level
Rollup}\label{stage-4-question-level-rollup}

\subsection{Objective}\label{objective-3}

Convert 1,598 pair-level results to 380 question-level consolidability
assessments with triage.

\subsection{Challenge: Pair-Level
Inflation}\label{challenge-pair-level-inflation}

Each of 380 source questions was compared against multiple ACS
questions: - \textbf{Average}: 4.2 comparisons per question -
\textbf{Range}: 1-12 comparisons per question - \textbf{Problem}: One F1
match makes a question consolidable, even if 11 other comparisons failed

\textbf{Solution}: Best-match rollup + two-axis triage.

\subsection{Best-Match Selection}\label{best-match-selection}

For each source question: 1. Identify all pairs involving that question
2. Select pair with highest consolidability: - F1 \textgreater{} F2
\textgreater{} F3 (feasibility hierarchy) - Within same feasibility,
highest Borda score 3. Output: One best match per question

\textbf{Result}: 380 question-level records with best ACS match.

\subsection{Scoring Methods (Ensemble
Approach)}\label{scoring-methods-ensemble-approach}

We computed four complementary scores for each pair:

\subsubsection{1. Composite Score
(Baseline)}\label{composite-score-baseline}

\begin{verbatim}
score = {
  1.0 if unanimous F1/F2,
  0.5 if majority F1/F2,
  0.0 if unanimous F3,
  NA otherwise
}
\end{verbatim}

\subsubsection{2. Entropy Score
(Stability)}\label{entropy-score-stability}

\begin{verbatim}
H = -Σ p(x) log₂ p(x)
normalized to [0, 1]
\end{verbatim}

\begin{itemize}
\tightlist
\item
  Low entropy = high agreement (stable)
\item
  High entropy = disagreement (unstable)
\end{itemize}

\subsubsection{3. Bayesian Score
(Probabilistic)}\label{bayesian-score-probabilistic}

\begin{verbatim}
P(consolidable | ratings) = (positives + α) / (total + α + β)
\end{verbatim}

\begin{itemize}
\tightlist
\item
  Incorporates prior beliefs
\item
  Shrinks toward population mean
\end{itemize}

\subsubsection{4. Borda Count (Direction)}\label{borda-count-direction}

\begin{verbatim}
Borda = Σ ranks / max_possible_ranks
\end{verbatim}

\begin{itemize}
\tightlist
\item
  F1 = 2 points, F2 = 1 point, F3 = 0 points
\item
  Captures direction of ensemble
\end{itemize}

\subsection{Two-Axis Triage Framework}\label{two-axis-triage-framework}

\textbf{Purpose}: Separate ``what's the answer?'' (Borda) from ``how
much did they argue?'' (Entropy).

\textbf{Thresholds}: Median split on question-level best-match scores -
Borda median: 0.167 - Entropy median: 0.330

\textbf{Why question-level medians?} - Pair-level distribution dominated
by unanimous F3 pairs (median ≈ 0) - Question-level reflects actual
decision space for triage - N=380 questions, not N=1,598 pairs

\textbf{Quadrant Definitions}:

\begin{longtable}[]{@{}
  >{\raggedright\arraybackslash}p{(\linewidth - 8\tabcolsep) * \real{0.2000}}
  >{\raggedright\arraybackslash}p{(\linewidth - 8\tabcolsep) * \real{0.1400}}
  >{\raggedright\arraybackslash}p{(\linewidth - 8\tabcolsep) * \real{0.1800}}
  >{\raggedright\arraybackslash}p{(\linewidth - 8\tabcolsep) * \real{0.3200}}
  >{\raggedright\arraybackslash}p{(\linewidth - 8\tabcolsep) * \real{0.1600}}@{}}
\toprule\noalign{}
\begin{minipage}[b]{\linewidth}\raggedright
Quadrant
\end{minipage} & \begin{minipage}[b]{\linewidth}\raggedright
Borda
\end{minipage} & \begin{minipage}[b]{\linewidth}\raggedright
Entropy
\end{minipage} & \begin{minipage}[b]{\linewidth}\raggedright
Interpretation
\end{minipage} & \begin{minipage}[b]{\linewidth}\raggedright
Action
\end{minipage} \\
\midrule\noalign{}
\endhead
\bottomrule\noalign{}
\endlastfoot
\textbf{Q1} & High & High & Confident consolidable & Auto-accept,
spot-check \\
\textbf{Q2} & Low & High & Confident non-consolidable & Auto-reject, low
priority \\
\textbf{Q3} & High & Low & Uncertain accept (leaning yes but contested)
& \textbf{Expert review priority} \\
\textbf{Q4} & Low & Low & Uncertain reject (genuinely ambiguous) &
Expert review secondary \\
\end{longtable}

\textbf{Rationale}: - Q1 (High Borda, High Entropy): Models agree it's
consolidable → trust - Q2 (Low Borda, High Entropy): Models agree it's
not consolidable → trust - Q3 (High Borda, Low Entropy): Leaning yes but
unstable → expert validate - Q4 (Low Borda, Low Entropy): Unclear
direction → expert clarify

\subsection{Triage Assignment Results}\label{triage-assignment-results}

\begin{longtable}[]{@{}llll@{}}
\toprule\noalign{}
Quadrant & Count & \% & Description \\
\midrule\noalign{}
\endhead
\bottomrule\noalign{}
\endlastfoot
Q1 & 151 & 39.7\% & Auto-accept (confident consolidable) \\
Q2 & 136 & 35.8\% & Auto-reject (confident non-consolidable) \\
Q3 & 40 & 10.5\% & Expert review (uncertain accept) \\
Q4 & 53 & 13.9\% & Expert review (uncertain reject) \\
\end{longtable}

\textbf{Expert review load}: 24.5\% (93 of 380 questions)

\begin{center}\rule{0.5\linewidth}{0.5pt}\end{center}

\section{Stage 5: Deliverables}\label{stage-5-deliverables}

\subsection{Objective}\label{objective-4}

Generate stakeholder-ready outputs for expert review and implementation.

\subsection{Expert Review Tables}\label{expert-review-tables}

Three CSV files generated: 1. \textbf{\texttt{expert\_review\_cps.csv}}
(240 CPS questions) 2. \textbf{\texttt{expert\_review\_foodaps.csv}}
(140 FoodAPS questions) 3.
\textbf{\texttt{expert\_review\_combined.csv}} (380 questions, all
surveys)

\textbf{Columns included}: - Source question text - Best ACS match text
- Feasibility (F1/F2/F3) - Barrier code (if F3) - Borda score
(direction) - Entropy score (stability) - Triage quadrant (Q1-Q4) -
Confidence flags - Reasoning summary

\subsection{Visualizations}\label{visualizations}

All figures saved to \texttt{output/visuals/}: - \textbf{Consolidation
rates} by survey (bar chart) - \textbf{Barrier distribution} (pie chart)
- \textbf{Expert review load} (stacked bar) - \textbf{Triage quadrant
heatmap} - \textbf{Harmonization code distribution} (pair-level) -
\textbf{Question consolidation distribution} (question-level by survey)

\subsection{Summary Statistics}\label{summary-statistics}

\texttt{output/analysis/stage4\_survey\_summary.json}:

\begin{Shaded}
\begin{Highlighting}[]
\FunctionTok{\{}
  \DataTypeTok{"CPS"}\FunctionTok{:} \FunctionTok{\{}
    \DataTypeTok{"total"}\FunctionTok{:} \DecValTok{240}\FunctionTok{,}
    \DataTypeTok{"consolidable"}\FunctionTok{:} \DecValTok{100}\FunctionTok{,}
    \DataTypeTok{"rate"}\FunctionTok{:} \FloatTok{0.417}
  \FunctionTok{\},}
  \DataTypeTok{"FoodAPS"}\FunctionTok{:} \FunctionTok{\{}
    \DataTypeTok{"total"}\FunctionTok{:} \DecValTok{140}\FunctionTok{,}
    \DataTypeTok{"consolidable"}\FunctionTok{:} \DecValTok{68}\FunctionTok{,}
    \DataTypeTok{"rate"}\FunctionTok{:} \FloatTok{0.486}
  \FunctionTok{\}}
\FunctionTok{\}}
\end{Highlighting}
\end{Shaded}

\begin{center}\rule{0.5\linewidth}{0.5pt}\end{center}

\section{Quality Assurance}\label{quality-assurance}

\subsection{Validation Checks}\label{validation-checks}

All pipeline stages include automated validation: 1. \textbf{Data
completeness}: No missing required fields 2. \textbf{Schema compliance}:
Output matches expected format 3. \textbf{Logical consistency}: No F1/F2
pairs with barrier codes 4. \textbf{Agreement metrics}: Kappa
\textgreater{} 0.75 threshold 5. \textbf{Rollup integrity}: Every
question has exactly one best match

\textbf{See}: \texttt{docs/ANALYSIS\_VV\_PLAN.md} for complete V\&V
checklist.

\subsection{Reproducibility}\label{reproducibility}

All analysis scripts are deterministic and produce identical results
when re-run: - Fixed random seeds - Checkpoint/resume capability -
Version-controlled prompts - Documented model versions

\begin{center}\rule{0.5\linewidth}{0.5pt}\end{center}

\section{Methodological Limitations}\label{methodological-limitations}

\subsection{1. Model Selection}\label{model-selection}

\begin{itemize}
\tightlist
\item
  Only three frontier models tested
\item
  Results may differ with other models
\item
  Behavioral profiles may change with model updates
\end{itemize}

\subsection{2. Prompt Engineering}\label{prompt-engineering}

\begin{itemize}
\tightlist
\item
  Structured prompts guide but don't guarantee consistency
\item
  Different prompt phrasings could yield different results
\item
  No systematic prompt optimization performed
\end{itemize}

\subsection{3. Barrier Taxonomy
Application}\label{barrier-taxonomy-application}

\begin{itemize}
\tightlist
\item
  Models apply human-designed taxonomy, not discovered patterns
\item
  Taxonomy completeness not validated empirically
\item
  Some pairs may have multiple applicable codes (we select primary)
\end{itemize}

\subsection{4. Triage Framework}\label{triage-framework}

\begin{itemize}
\tightlist
\item
  Two-axis approach is operational heuristic, not theoretical optimum
\item
  Threshold selection (median split) is pragmatic, not optimized
\item
  Quadrant interpretation requires domain context
\end{itemize}

\subsection{5. Scope}\label{scope-1}

\begin{itemize}
\tightlist
\item
  Only CPS and FoodAPS analyzed (not full federal survey ecosystem)
\item
  Only ACS as target (other surveys could serve as consolidation
  targets)
\item
  English-language surveys only
\end{itemize}

\begin{center}\rule{0.5\linewidth}{0.5pt}\end{center}

\section{Next Steps After This
Report}\label{next-steps-after-this-report}

\begin{enumerate}
\def\labelenumi{\arabic{enumi}.}
\tightlist
\item
  \textbf{Expert validation}: Subject-matter experts review Q3/Q4
  questions
\item
  \textbf{Consolidation pilots}: Test F1 recommendations in practice
\item
  \textbf{Expansion}: Apply methodology to additional survey pairs
\item
  \textbf{Refinement}: Incorporate expert feedback to improve
  classification
\end{enumerate}

\textbf{Key principle}: AI assists, experts decide. This methodology
accelerates analysis but does not replace human judgment.

\bookmarksetup{startatroot}

\chapter{Results}\label{results}

\section{Consolidation Rates}\label{consolidation-rates}

\subsection{Overall Question-Level
Results}\label{overall-question-level-results}

\begin{figure}[H]

{\centering \pandocbounded{\includegraphics[keepaspectratio]{sections/figures/consolidation_rates.png}}

}

\caption{Consolidation Rates by Survey}

\end{figure}%

\begin{longtable}[]{@{}lllllll@{}}
\toprule\noalign{}
Survey & Total Questions & F1 & F2 & F3 & Consolidable & Rate \\
\midrule\noalign{}
\endhead
\bottomrule\noalign{}
\endlastfoot
\textbf{CPS} & 240 & 37 & 63 & 140 & 100 & \textbf{41.7\%} \\
\textbf{FoodAPS} & 140 & 23 & 45 & 72 & 68 & \textbf{48.6\%} \\
\textbf{Combined} & 380 & 60 & 108 & 212 & 168 & \textbf{44.2\%} \\
\end{longtable}

\textbf{Key Finding}: Nearly half of surveyed questions have at least
one consolidable ACS match.

\subsection{Consolidation Type
Breakdown}\label{consolidation-type-breakdown}

\begin{itemize}
\tightlist
\item
  \textbf{15.8\%} (60/380) are F1 - directly consolidable with simple
  recoding
\item
  \textbf{28.4\%} (108/380) are F2 - consolidable with statistical
  adjustment
\item
  \textbf{55.8\%} (212/380) are F3 - fundamentally incompatible
\end{itemize}

\textbf{Interpretation}: - F1 questions are ``low-hanging fruit'' -
immediate consolidation candidates - F2 questions require modeling but
remain viable - F3 questions need deeper examination of barriers

\begin{center}\rule{0.5\linewidth}{0.5pt}\end{center}

\section{Barrier Analysis}\label{barrier-analysis}

\subsection{Pair-Level Distribution}\label{pair-level-distribution}

\begin{figure}[H]

{\centering \pandocbounded{\includegraphics[keepaspectratio]{sections/figures/barrier_distribution.png}}

}

\caption{Barrier Distribution (1,598 pairs)}

\end{figure}%

Out of 1,283 F3 pairs (incompatible):

\begin{longtable}[]{@{}
  >{\raggedright\arraybackslash}p{(\linewidth - 6\tabcolsep) * \real{0.3696}}
  >{\raggedright\arraybackslash}p{(\linewidth - 6\tabcolsep) * \real{0.1522}}
  >{\raggedright\arraybackslash}p{(\linewidth - 6\tabcolsep) * \real{0.1957}}
  >{\raggedright\arraybackslash}p{(\linewidth - 6\tabcolsep) * \real{0.2826}}@{}}
\toprule\noalign{}
\begin{minipage}[b]{\linewidth}\raggedright
Barrier Category
\end{minipage} & \begin{minipage}[b]{\linewidth}\raggedright
Count
\end{minipage} & \begin{minipage}[b]{\linewidth}\raggedright
\% of F3
\end{minipage} & \begin{minipage}[b]{\linewidth}\raggedright
Description
\end{minipage} \\
\midrule\noalign{}
\endhead
\bottomrule\noalign{}
\endlastfoot
\textbf{CC} (Construct/Concept) & 1,238 & \textbf{96.5\%} & Fundamental
concept differences \\
\textbf{TC} (Temporal) & 24 & 1.9\% & Reference period or timing
issues \\
\textbf{RS} (Response Scale) & 13 & 1.0\% & Scale type or category
differences \\
\textbf{PC} (Population) & 6 & 0.5\% & Universe or coverage
differences \\
\textbf{MC} (Mode/Context) & 1 & 0.1\% & Interview mode or routing \\
\textbf{PM} (Processing) & 1 & 0.1\% & Coding or metadata issues \\
\end{longtable}

\textbf{Key Finding}: 97\% of failures stem from construct/concept (CC)
barriers, not operational issues.

\subsection{Construct Barrier
Sub-Codes}\label{construct-barrier-sub-codes}

Among CC barriers (1,238 pairs):

\begin{longtable}[]{@{}lllll@{}}
\toprule\noalign{}
Sub-Code & Count & \% of CC & \% of F3 & Description \\
\midrule\noalign{}
\endhead
\bottomrule\noalign{}
\endlastfoot
\textbf{CC.1} & 900 & 72.7\% & 70.1\% & Concept definition
differences \\
\textbf{CC.2} & 250 & 20.2\% & 19.5\% & Operationalization
differences \\
\textbf{CC.4} & 88 & 7.1\% & 6.9\% & Scope inclusion differences \\
\end{longtable}

\textbf{Interpretation}: - \textbf{CC.1 dominance}: Most incompatible
pairs measure fundamentally different constructs - \textbf{Not fixable}:
These are not measurement errors or survey design choices - they reflect
different research goals - \textbf{Example}: ``Do you want to work
full-time?'' (intention) vs ``Did you work last week?'' (behavior)

\subsection{Question-Level Barrier
Distribution}\label{question-level-barrier-distribution}

\begin{figure}[H]

{\centering \pandocbounded{\includegraphics[keepaspectratio]{sections/figures/question_consolidation_distribution.png}}

}

\caption{Question Consolidation by Survey}

\end{figure}%

At the question level (212 F3 questions):

\begin{longtable}[]{@{}lll@{}}
\toprule\noalign{}
Barrier & Count & \% of F3 Questions \\
\midrule\noalign{}
\endhead
\bottomrule\noalign{}
\endlastfoot
\textbf{CC.1} & 163 & \textbf{76.9\%} \\
\textbf{CC.2} & 38 & 17.9\% \\
\textbf{CC.4} & 4 & 1.9\% \\
\textbf{Other} & 7 & 3.3\% \\
\end{longtable}

\textbf{Key Finding}: CC.1 concentration is even higher at question
level (76.9\% vs 70.1\%), indicating that questions without any
consolidable path tend to have fundamental concept mismatches.

\begin{center}\rule{0.5\linewidth}{0.5pt}\end{center}

\section{Agreement Statistics}\label{agreement-statistics}

\subsection{Inter-Rater Reliability}\label{inter-rater-reliability}

\begin{figure}[H]

{\centering \includegraphics[width=0.7\linewidth,height=\textheight,keepaspectratio]{sections/figures/agreement_heatmap.png}

}

\caption{Agreement Heatmap}

\end{figure}%

\begin{longtable}[]{@{}lll@{}}
\toprule\noalign{}
Model Pair & Cohen's κ & Interpretation \\
\midrule\noalign{}
\endhead
\bottomrule\noalign{}
\endlastfoot
\textbf{OpenAI - Anthropic} & 0.845 & Almost perfect \\
\textbf{OpenAI - Google} & 0.796 & Substantial \\
\textbf{Anthropic - Google} & 0.833 & Almost perfect \\
\end{longtable}

\textbf{Fleiss' κ (three-way)}: 0.833 (almost perfect)

\subsection{Agreement by Classification
Level}\label{agreement-by-classification-level}

\begin{longtable}[]{@{}
  >{\raggedright\arraybackslash}p{(\linewidth - 4\tabcolsep) * \real{0.2000}}
  >{\raggedright\arraybackslash}p{(\linewidth - 4\tabcolsep) * \real{0.4286}}
  >{\raggedright\arraybackslash}p{(\linewidth - 4\tabcolsep) * \real{0.3714}}@{}}
\toprule\noalign{}
\begin{minipage}[b]{\linewidth}\raggedright
Level
\end{minipage} & \begin{minipage}[b]{\linewidth}\raggedright
Agreement Rate
\end{minipage} & \begin{minipage}[b]{\linewidth}\raggedright
Description
\end{minipage} \\
\midrule\noalign{}
\endhead
\bottomrule\noalign{}
\endlastfoot
\textbf{Feasibility (F1/F2/F3)} & 84.5\% & High-level harmonization
potential \\
\textbf{Barrier L1 (CC/TC/RS)} & 91.2\% & Barrier category \\
\textbf{Barrier L2 (CC.1/CC.2)} & 78.9\% & Specific barrier sub-code \\
\end{longtable}

\textbf{Interpretation}: Models show strong agreement on broad
feasibility judgments, with expected variation on granular barrier
codes.

\subsection{Confidence Correlation}\label{confidence-correlation}

\begin{longtable}[]{@{}lll@{}}
\toprule\noalign{}
Confidence Level & Pair Count & Agreement Rate \\
\midrule\noalign{}
\endhead
\bottomrule\noalign{}
\endlastfoot
\textbf{HIGH} & 1,089 (68.1\%) & 91.3\% \\
\textbf{MODERATE} & 389 (24.3\%) & 76.8\% \\
\textbf{LOW} & 120 (7.5\%) & 54.2\% \\
\end{longtable}

\textbf{Finding}: Model-reported confidence predicts agreement.
High-confidence pairs show 91\% agreement, validating confidence as a
triage signal.

\begin{center}\rule{0.5\linewidth}{0.5pt}\end{center}

\section{Expert Review Load}\label{expert-review-load}

\begin{figure}[H]

{\centering \pandocbounded{\includegraphics[keepaspectratio]{sections/figures/expert_review_load.png}}

}

\caption{Expert Review Load Distribution}

\end{figure}%

\subsection{Triage Quadrant
Distribution}\label{triage-quadrant-distribution}

\begin{longtable}[]{@{}
  >{\raggedright\arraybackslash}p{(\linewidth - 8\tabcolsep) * \real{0.2500}}
  >{\raggedright\arraybackslash}p{(\linewidth - 8\tabcolsep) * \real{0.1750}}
  >{\raggedright\arraybackslash}p{(\linewidth - 8\tabcolsep) * \real{0.0750}}
  >{\raggedright\arraybackslash}p{(\linewidth - 8\tabcolsep) * \real{0.3000}}
  >{\raggedright\arraybackslash}p{(\linewidth - 8\tabcolsep) * \real{0.2000}}@{}}
\toprule\noalign{}
\begin{minipage}[b]{\linewidth}\raggedright
Quadrant
\end{minipage} & \begin{minipage}[b]{\linewidth}\raggedright
Count
\end{minipage} & \begin{minipage}[b]{\linewidth}\raggedright
\%
\end{minipage} & \begin{minipage}[b]{\linewidth}\raggedright
Confidence
\end{minipage} & \begin{minipage}[b]{\linewidth}\raggedright
Action
\end{minipage} \\
\midrule\noalign{}
\endhead
\bottomrule\noalign{}
\endlastfoot
\textbf{Q1} (High Borda, High Entropy) & 151 & 39.7\% & Confident
consolidable & Auto-accept \\
\textbf{Q2} (Low Borda, High Entropy) & 136 & 35.8\% & Confident
non-consolidable & Auto-reject \\
\textbf{Q3} (High Borda, Low Entropy) & 40 & 10.5\% & Uncertain accept &
\textbf{Expert priority} \\
\textbf{Q4} (Low Borda, Low Entropy) & 53 & 13.9\% & Uncertain reject &
Expert secondary \\
\end{longtable}

\textbf{Expert Review Load}: 93 questions (24.5\%) flagged for human
review

\subsection{Load Reduction}\label{load-reduction}

Traditional approach: \textbf{380 questions × 100\%} = 380 expert
reviews

AI-assisted approach: - \textbf{287 auto-processed} (75.5\%) -
spot-check only - \textbf{93 expert reviews} (24.5\%) - full human
judgment

\textbf{Efficiency gain}: 75.5\% reduction in expert review workload,
focusing effort on genuinely uncertain cases.

\begin{center}\rule{0.5\linewidth}{0.5pt}\end{center}

\section{Example Cases}\label{example-cases}

\subsection{High Consolidability (F1)}\label{high-consolidability-f1}

\textbf{Source (CPS)}: ``How many hours per week do you USUALLY work at
your main job?''

\textbf{Target (ACS)}: ``What is your best estimate of the number of
hours per week you usually work at this rate?''

\textbf{Verdict}: F1 (directly consolidable) - Same construct (usual
work hours) - Same reference period (typical week) - Same measurement
approach - Minor wording differences immaterial

\textbf{Scores}: Borda=1.0, Entropy=1.0 → Q1 (confident consolidable)

\begin{center}\rule{0.5\linewidth}{0.5pt}\end{center}

\subsection{Medium Consolidability
(F2)}\label{medium-consolidability-f2}

\textbf{Source (CPS)}: ``EXCLUDING overtime, tips, commissions - what is
your hourly rate of pay?''

\textbf{Target (ACS)}: ``What is your hourly rate of pay?''

\textbf{Verdict}: F2 (consolidable with transformation) - Same construct
(hourly wage) - Different scope (CPS excludes components, ACS includes
all) - Statistical adjustment needed: Scale CPS response to account for
excluded components - Requires bridging study or imputation model

\textbf{Barrier}: CC.4 (scope inclusion differences)

\textbf{Scores}: Borda=0.5, Entropy=0.67 → Q1 (confident consolidable
with adjustment)

\begin{center}\rule{0.5\linewidth}{0.5pt}\end{center}

\subsection{Not Consolidable (F3)}\label{not-consolidable-f3}

\textbf{Source (CPS)}: ``Do you WANT to work full-time (35+ hours per
week)?''

\textbf{Target (ACS)}: ``Did this person WORK for pay last week?''

\textbf{Verdict}: F3 (incompatible) - Different constructs: Intention vs
behavior - CPS measures labor force preferences - ACS measures actual
work status - No statistical transformation can bridge this conceptual
gap

\textbf{Barrier}: CC.1 (concept definition differences)

\textbf{Scores}: Borda=0.0, Entropy=0.33 → Q2 (confident
non-consolidable)

\begin{center}\rule{0.5\linewidth}{0.5pt}\end{center}

\section{Topic-Level Breakdown}\label{topic-level-breakdown}

\subsection{Consolidability by Question
Topic}\label{consolidability-by-question-topic}

\begin{longtable}[]{@{}lllll@{}}
\toprule\noalign{}
Topic & Total & Consolidable & Rate & Notes \\
\midrule\noalign{}
\endhead
\bottomrule\noalign{}
\endlastfoot
Demographics & 45 & 38 & 84.4\% & High overlap (age, sex, race) \\
Economics & 128 & 62 & 48.4\% & Moderate (income definitions vary) \\
Employment & 147 & 52 & 35.4\% & Low (CPS specialization) \\
Health & 32 & 12 & 37.5\% & Low (different health constructs) \\
Social & 28 & 4 & 14.3\% & Very low (specialized topics) \\
\end{longtable}

\textbf{Interpretation}: - Demographics have high consolidation
potential (expected - standard measures) - Employment shows low
consolidation despite CPS-ACS overlap (CPS is specialized) - Social
topics have very low consolidation (specialized constructs)

\begin{center}\rule{0.5\linewidth}{0.5pt}\end{center}

\section{Survey-Specific Insights}\label{survey-specific-insights}

\subsection{CPS (Current Population
Survey)}\label{cps-current-population-survey}

\begin{itemize}
\tightlist
\item
  \textbf{Focus}: Labor force statistics
\item
  \textbf{Consolidation rate}: 41.7\% (100/240)
\item
  \textbf{Challenge}: Many employment questions are CPS-specific (job
  search, work preferences, multiple job holders)
\item
  \textbf{Opportunity}: Demographics and basic employment status
  consolidable
\end{itemize}

\subsection{FoodAPS (Food Acquisition and Purchase
Survey)}\label{foodaps-food-acquisition-and-purchase-survey}

\begin{itemize}
\tightlist
\item
  \textbf{Focus}: Food security and purchasing behavior
\item
  \textbf{Consolidation rate}: 48.6\% (68/140)
\item
  \textbf{Challenge}: Specialized food security batteries not in ACS
\item
  \textbf{Opportunity}: Demographics and household composition
  consolidable
\end{itemize}

\textbf{Why FoodAPS higher?} - Shorter questionnaire (140 vs 240
questions) - More demographic questions relative to specialized content
- Less employment specialization than CPS

\begin{center}\rule{0.5\linewidth}{0.5pt}\end{center}

\section{Summary Statistics}\label{summary-statistics-1}

\subsection{At a Glance}\label{at-a-glance}

\begin{verbatim}
Input: 380 source questions (240 CPS + 140 FoodAPS)
Pairs analyzed: 1,598 (avg 4.2 comparisons per question)

Output:
- 168 consolidable questions (44.2%)
  - 60 F1 (15.8%) - direct recode
  - 108 F2 (28.4%) - statistical adjustment
- 212 non-consolidable (55.8%)
  - 97% due to CC barriers
  - 70% specifically CC.1

Expert review: 93 questions (24.5%)
Auto-processed: 287 questions (75.5%)

Inter-rater agreement: κ = 0.845 (almost perfect)
\end{verbatim}

\begin{center}\rule{0.5\linewidth}{0.5pt}\end{center}

\textbf{Next chapter}: Interpretation of these findings and implications
for federal survey harmonization practice.

\bookmarksetup{startatroot}

\chapter{Discussion}\label{discussion}

\section{Interpretation of Findings}\label{interpretation-of-findings}

\subsection{The 44\% Consolidation
Finding}\label{the-44-consolidation-finding}

Our analysis reveals that 44.2\% of questions (168 of 380) have at least
one consolidable ACS match. This finding must be interpreted carefully:

\textbf{What this means:} - Substantial overlap exists between federal
surveys - Nearly half of questions measure concepts also captured in ACS
- Consolidation is technically feasible for these questions

\textbf{What this does NOT mean:} - These questions should be eliminated
from source surveys - Consolidation is always desirable or appropriate -
Survey-specific context and use cases are unimportant

\textbf{Policy vs.~Technical Distinction}: This analysis establishes
\emph{technical feasibility} of consolidation. The decision to actually
consolidate questions requires stakeholder input, consideration of
survey-specific research goals, and assessment of data quality
trade-offs.

\subsection{The CC Barrier Dominance
(97\%)}\label{the-cc-barrier-dominance-97}

The finding that 97\% of incompatible pairs fail due to
Construct/Concept (CC) barriers, with 70\% specifically CC.1 (concept
definition differences), has important implications:

\subsubsection{This is NOT a Fixable
Problem}\label{this-is-not-a-fixable-problem}

Unlike operational barriers (mode effects, scale differences, temporal
misalignment), construct differences cannot be resolved through: -
Survey redesign - Question rewording - Statistical adjustment -
Post-processing harmonization

\textbf{Example}: A question asking ``Do you want to work full-time?''
measures \emph{labor force preferences}, while ``Did you work last
week?'' measures \emph{actual employment status}. No amount of
methodological sophistication can make these conceptually equivalent.

\subsubsection{This Validates Survey
Specialization}\label{this-validates-survey-specialization}

The CC.1 dominance suggests that federal surveys serve distinct research
purposes: - \textbf{CPS}: Detailed employment dynamics (job search,
multiple jobs, work preferences) - \textbf{FoodAPS}: Food security
measurement (food purchasing, SNAP utilization, food insecurity
batteries) - \textbf{ACS}: Broad demographic and economic snapshot

These specializations are intentional and valuable. Consolidation makes
sense for overlapping demographic measures, but each survey must retain
its unique substantive content.

\subsection{AI-Assisted Methods
Performance}\label{ai-assisted-methods-performance}

\subsubsection{High Agreement (κ =
0.845)}\label{high-agreement-ux3ba-0.845}

The strong inter-rater reliability across three independent LLMs
validates that: 1. Harmonization judgments are consistent and
reproducible 2. Different model architectures converge on similar
classifications 3. The framework (F1/F2/F3 + barrier codes) is
well-specified enough for models to apply consistently

\subsubsection{Expert Review Load Reduction
(75\%)}\label{expert-review-load-reduction-75}

By auto-processing 287 of 380 questions (75.5\%), the AI-assisted
approach achieves: - \textbf{Time savings}: Days instead of weeks for
initial classification - \textbf{Expert focus}: Human effort
concentrated on 93 genuinely uncertain cases - \textbf{Scalability}:
Methodology can be applied to additional survey pairs

\textbf{Important}: ``Auto-processing'' means high-confidence
classifications, not zero human oversight. All results should undergo
expert spot-checking and validation.

\subsection{The Two-Axis Triage
Framework}\label{the-two-axis-triage-framework}

Our Borda-Entropy triage approach represents an operational contribution
rather than a theoretical advance:

\textbf{What it is}: - Pragmatic heuristic for separating ``what's the
answer?'' (Borda) from ``how much did they argue?'' (Entropy) -
Operationalization of ensemble uncertainty for expert review
prioritization - Median-split thresholding for tractable quadrant
assignment

\textbf{What it is NOT}: - Novel theoretical framework requiring
extensive citation - Optimized decision boundary (we used pragmatic
median split) - Replacement for subject-matter expertise

\textbf{Framing}: This is a useful operational tool that performed well
in this context. We document it clearly for replication but do not claim
theoretical novelty. The framing in
\texttt{docs/stage4\_ensemble\_methodology.md} explicitly treats this as
operational methodology, not research contribution.

\begin{center}\rule{0.5\linewidth}{0.5pt}\end{center}

\section{Limitations}\label{limitations}

\subsection{1. Survey Coverage}\label{survey-coverage}

\textbf{Limitation}: Only two source surveys (CPS and FoodAPS) analyzed
against one target (ACS).

\textbf{Implications}: - Findings may not generalize to other federal
surveys - ACS-centric perspective (other surveys could serve as
consolidation targets) - Missing surveys with different topical focus
(health, housing, education)

\textbf{Mitigation}: Results demonstrate proof-of-concept; expansion to
additional surveys is straightforward using same methodology.

\subsection{2. Model Selection and
Behavior}\label{model-selection-and-behavior}

\textbf{Limitation}: Three frontier LLMs tested, each with documented
behavioral profiles: - OpenAI: Slight optimism toward F2 - Anthropic:
Balanced, detailed reasoning - Google: Conservative, deferential

\textbf{Implications}: - Results may differ with other models -
Behavioral profiles may change with model updates - No systematic
exploration of prompt engineering

\textbf{Mitigation}: High inter-rater agreement (κ = 0.845) suggests
results are robust across models. Arbitration pattern further reduces
single-model bias.

\subsection{3. Google Rate Limiting}\label{google-rate-limiting}

\textbf{Limitation}: Google API rate-limited at 751 pairs (47\% of
dataset).

\textbf{Implications}: - Three-way agreement analysis limited to 751
pairs - Primary findings based on OpenAI-Anthropic two-way comparison -
Google results used for validation subset only

\textbf{Mitigation}: Two-way agreement (κ = 0.845) is sufficient for
validation. Subset with Google (κ = 0.833) confirms consistency.

\subsection{4. Pair-Level Inflation}\label{pair-level-inflation}

\textbf{Limitation}: Analyzing 1,598 pairs for 380 questions creates
statistical dependencies.

\textbf{Challenge}: - Not all pairs are independent observations -
Multiple comparisons per question inflates denominator - Naive
pair-level statistics misleading

\textbf{Mitigation}: Question-level rollup (Stage 4) addresses this by
selecting one best match per question. Primary findings reported at
question level (N=380), not pair level (N=1,598).

\subsection{5. Barrier Taxonomy
Application}\label{barrier-taxonomy-application-1}

\textbf{Limitation}: Models apply human-designed taxonomy
(DataSHaPER/Maelstrom framework), not discovered patterns.

\textbf{Implications}: - Taxonomy completeness not empirically validated
- Some pairs may have multiple applicable barriers (we select primary) -
Granular sub-codes (e.g., CC.1 vs CC.2) subject to interpretation

\textbf{Mitigation}: High inter-rater agreement on barrier categories
(L1: 91.2\%) suggests taxonomy is well-specified. Expert validation will
further refine barrier assignments.

\subsection{6. Validation Status}\label{validation-status}

\textbf{Limitation}: Expert validation of classifications is pending.

\textbf{Implications}: - Findings represent AI ensemble judgments, not
validated ground truth - Classification errors possible, especially in
uncertain cases - Consolidation recommendations require stakeholder
review

\textbf{Mitigation}: Triage framework explicitly routes uncertain cases
(Q3/Q4) to expert review. High-confidence auto-processed cases (Q1/Q2)
recommended for spot-checking.

\begin{center}\rule{0.5\linewidth}{0.5pt}\end{center}

\section{Methodological
Contributions}\label{methodological-contributions}

\subsection{1. Operational Framework for AI-Assisted
Harmonization}\label{operational-framework-for-ai-assisted-harmonization}

This work demonstrates that multi-model ensemble with arbitration can: -
Classify survey question pairs consistently (κ = 0.845) - Identify
consolidation candidates at scale (1,598 pairs in days) - Route
uncertain cases to experts (75\% auto-processed)

\textbf{Contribution}: Proof-of-concept that AI-assisted methods can
accelerate survey harmonization without replacing expert judgment.

\subsection{2. Two-Axis Triage for Expert
Review}\label{two-axis-triage-for-expert-review}

The Borda-Entropy approach provides: - Separation of direction (what's
the answer?) from stability (did they agree?) - Quadrant-based
prioritization (Q3 priority, Q4 secondary) - Quantifiable expert review
load (24.5\%)

\textbf{Contribution}: Operational heuristic for prioritizing expert
effort. Not a theoretical advance, but a useful tool.

\subsection{3. Question-Level Rollup Addressing Pair
Inflation}\label{question-level-rollup-addressing-pair-inflation}

Converting pair-level classifications to question-level consolidability
via best-match rollup: - Resolves statistical dependency issues -
Provides actionable stakeholder deliverable (one recommendation per
question) - Enables question-level triage

\textbf{Contribution}: Methodological solution to multi-comparison
problem in pairwise analysis.

\subsection{4. Empirical Application of DataSHaPER Framework to Federal
Surveys}\label{empirical-application-of-datashaper-framework-to-federal-surveys}

First systematic application of DataSHaPER/Maelstrom harmonization
framework to U.S. federal survey ecosystem: - Quantifies consolidation
potential (\textasciitilde44\%) - Identifies barrier distribution (97\%
CC) - Validates framework applicability beyond epidemiology

\textbf{Contribution}: Demonstrates that survey harmonization frameworks
from international and health research translate to federal statistical
surveys.

\begin{center}\rule{0.5\linewidth}{0.5pt}\end{center}

\section{Implications for Practice}\label{implications-for-practice}

\subsection{For Federal Statistical
Agencies}\label{for-federal-statistical-agencies}

\textbf{Finding}: 44\% of questions consolidable suggests meaningful
burden reduction is possible.

\textbf{Actionable Steps}: 1. \textbf{Prioritize F1 questions} (60
questions, 15.8\%) as ``low-hanging fruit'' - these can be consolidated
with minimal effort 2. \textbf{Evaluate F2 questions} (108 questions,
28.4\%) case-by-case - assess whether statistical adjustment is worth
the effort 3. \textbf{Respect F3 questions} (212 questions, 55.8\%) -
these serve distinct research purposes and should remain survey-specific

\textbf{Important}: Consolidation is a tool for burden reduction, not a
mandate. Survey-specific research goals must guide decisions.

\subsection{For Survey Methodologists}\label{for-survey-methodologists}

\textbf{Finding}: AI-assisted methods reduce expert review load by 75\%.

\textbf{Implications}: - \textbf{Triage, don't replace}: Use AI for
initial classification, reserve experts for uncertain cases -
\textbf{Validate liberally}: Spot-check auto-processed classifications
to ensure quality - \textbf{Iterate}: As experts review, refine prompts
and threshold based on feedback

\textbf{Workflow}:

\begin{verbatim}
AI Ensemble → High Confidence (Auto-process) + Low Confidence (Expert Review)
              ↓                                ↓
         Spot-check validation          Full expert judgment
              ↓                                ↓
         Final classification          Refine methodology
\end{verbatim}

\subsection{For Data Users}\label{for-data-users}

\textbf{Finding}: Consolidated questions enable cross-survey data
integration.

\textbf{Opportunities}: - Link CPS employment data with ACS housing data
using harmonized variables - Combine FoodAPS food security measures with
ACS income data - Perform longitudinal analysis across surveys with
consistent measures

\textbf{Caution}: Even F1 questions may have subtle differences
(sampling, mode, context). Users should assess fitness-for-purpose
before integrating data.

\begin{center}\rule{0.5\linewidth}{0.5pt}\end{center}

\section{Relation to Prior Work}\label{relation-to-prior-work}

\subsection{Reports 01-02}\label{reports-01-02}

This analysis builds on: - \textbf{Report 01}: Concept classification
established question taxonomy - \textbf{Report 02}: Pairwise comparison
identified 1,598 pairs for analysis

\textbf{Integration}: Each report's validated outputs feed forward,
demonstrating cumulative progress.

\subsection{Survey Harmonization
Literature}\label{survey-harmonization-literature}

Our findings align with prior work: - \textbf{Fortier et al.~(2011)}:
Found 38\% harmonization rate across 53 epidemiological studies - our
44\% is comparable - \textbf{Wolf et al.~(2016)}: Emphasized construct
barriers as primary challenge - validated by our 97\% CC finding

\textbf{Extension}: We apply established frameworks to federal surveys
at scale using AI assistance - first such application documented.

\begin{center}\rule{0.5\linewidth}{0.5pt}\end{center}

\section{Alternative Explanations}\label{alternative-explanations}

\subsection{Why 44\% and not Higher?}\label{why-44-and-not-higher}

One might expect higher consolidation rates given that federal surveys
all cover demographic and economic topics.

\textbf{Explanation}: 1. \textbf{Survey specialization}: CPS focuses on
employment, FoodAPS on food security - specialized questions don't
overlap with ACS general-purpose content 2. \textbf{Research goals
differ}: Surveys measure concepts at different granularities and from
different perspectives 3. \textbf{Policy-driven content}: Questions
often driven by specific policy needs (e.g., SNAP eligibility) not
present in other surveys

\textbf{Counterpoint}: 44\% is actually quite high given survey
specialization. Demographics show 84\% consolidation, validating that
overlap exists where expected.

\subsection{Why CC.1 Dominates?}\label{why-cc.1-dominates}

The 70\% CC.1 concentration (concept definition differences) might seem
surprising.

\textbf{Explanation}: 1. \textbf{Fundamental barriers filter early}:
Questions with fixable barriers (TC, RS, PC) often have alternative
matches that succeed 2. \textbf{Best-match selection}: We select highest
consolidability per question - minor barriers appear less frequently
because alternatives exist 3. \textbf{Survey design}: Federal surveys
intentionally measure different constructs, so CC barriers are the true
bottleneck

\textbf{Evidence}: Question-level shows even higher CC.1 (76.9\%),
confirming that questions without any consolidable path have fundamental
concept mismatches.

\begin{center}\rule{0.5\linewidth}{0.5pt}\end{center}

\section{Future Directions}\label{future-directions}

\subsection{Immediate Next Steps}\label{immediate-next-steps}

\begin{enumerate}
\def\labelenumi{\arabic{enumi}.}
\tightlist
\item
  \textbf{Expert validation}: Subject-matter experts review 93 flagged
  questions (Q3/Q4)
\item
  \textbf{Consolidation pilots}: Test F1 recommendations in practice
\item
  \textbf{Feedback incorporation}: Refine classifications based on
  expert input
\end{enumerate}

\subsection{Short-Term Extensions}\label{short-term-extensions}

\begin{enumerate}
\def\labelenumi{\arabic{enumi}.}
\tightlist
\item
  \textbf{Additional surveys}: Apply methodology to SIPP, NHIS, AHS
\item
  \textbf{Alternative targets}: Use CPS or FoodAPS as consolidation
  targets (not just ACS)
\item
  \textbf{Bidirectional analysis}: Identify ACS questions that could
  adopt specialized survey measures
\end{enumerate}

\subsection{Long-Term Applications}\label{long-term-applications}

\begin{enumerate}
\def\labelenumi{\arabic{enumi}.}
\tightlist
\item
  \textbf{Survey editing}: Use harmonization patterns to inform
  questionnaire design
\item
  \textbf{Imputation frameworks}: Leverage consolidable questions for
  cross-survey imputation
\item
  \textbf{Respondent burden modeling}: Quantify burden reduction from
  consolidation
\end{enumerate}

\begin{center}\rule{0.5\linewidth}{0.5pt}\end{center}

\textbf{Key Takeaway}: This analysis provides an evidence base for
federal survey harmonization decisions. The methodology is reproducible,
the findings are actionable, and the limitations are documented.
Stakeholders now have data-driven recommendations for which questions to
consolidate, which to harmonize with adjustment, and which to maintain
as survey-specific.

\bookmarksetup{startatroot}

\chapter{Conclusion}\label{conclusion}

\section{Summary of Findings}\label{summary-of-findings}

This report presents the first systematic analysis of federal survey
question consolidation potential using AI-assisted harmonization
methods. Analyzing 380 questions (240 CPS + 140 FoodAPS) against the
American Community Survey, we identified consolidation opportunities and
characterized barriers at scale.

\subsection{Key Results}\label{key-results-1}

\textbf{Consolidation Potential}: - \textbf{44.2\%} of questions have at
least one consolidable ACS match - 15.8\% (60 questions) are F1 -
directly consolidable - 28.4\% (108 questions) are F2 - consolidable
with statistical adjustment - 55.8\% (212 questions) are F3 -
fundamentally incompatible

\textbf{Barrier Analysis}: - \textbf{97\% of failures} stem from
Construct/Concept (CC) barriers - \textbf{70\%} specifically CC.1
(concept definition differences) - These are not fixable through survey
redesign or statistical methods

\textbf{AI Performance}: - \textbf{κ = 0.845} inter-rater agreement
(almost perfect) - \textbf{75\%} of questions auto-processed with high
confidence - \textbf{24.5\%} (93 questions) routed to expert review

\subsection{Implications}\label{implications-1}

\textbf{For Policy}: Substantial consolidation potential exists, but
decisions must balance burden reduction with survey-specific research
goals.

\textbf{For Methodology}: AI-assisted ensemble methods can accelerate
survey harmonization while preserving expert oversight.

\textbf{For Practice}: The two-axis triage framework (Borda direction ×
Entropy stability) provides an operational tool for prioritizing expert
review.

\begin{center}\rule{0.5\linewidth}{0.5pt}\end{center}

\section{Contributions}\label{contributions}

This work contributes:

\subsection{1. Empirical Evidence}\label{empirical-evidence}

First quantification of consolidation potential across federal surveys
using standardized harmonization framework. The 44\% finding provides
data for burden reduction planning.

\subsection{2. Barrier Characterization}\label{barrier-characterization}

Systematic application of DataSHaPER/Maelstrom barrier taxonomy
identifies that construct differences (not operational issues) prevent
consolidation. This clarifies where effort should focus.

\subsection{3. Operational Methodology}\label{operational-methodology}

Reproducible AI-assisted pipeline with: - Multi-model ensemble to reduce
bias - Arbitration for disagreement resolution - Two-axis triage for
expert review routing - Question-level rollup to address pair inflation

\subsection{4. Stakeholder Deliverables}\label{stakeholder-deliverables}

Expert review tables ready for validation and implementation, with 168
specific consolidation recommendations.

\begin{center}\rule{0.5\linewidth}{0.5pt}\end{center}

\section{Limitations Revisited}\label{limitations-revisited}

We acknowledge: - \textbf{Limited survey coverage}: Only CPS and FoodAPS
analyzed - \textbf{Model-specific results}: Based on three frontier LLMs
(OpenAI, Anthropic, Google) - \textbf{Pending validation}: Expert review
of classifications underway - \textbf{Pair inflation}: Addressed via
question-level rollup but inherent to exhaustive comparison

These limitations do not invalidate findings but contextualize their
generalizability.

\begin{center}\rule{0.5\linewidth}{0.5pt}\end{center}

\section{Future Work}\label{future-work}

\subsection{Immediate Priorities}\label{immediate-priorities}

\textbf{1. Expert Validation (Q1 2026)}

Subject-matter experts review: - \textbf{Q3 questions} (40 questions,
high priority): Leaning consolidable but contested - \textbf{Q4
questions} (53 questions, secondary): Genuinely ambiguous

\textbf{Outcome}: Refined classifications, validated recommendations

\textbf{2. Consolidation Pilots (Q2 2026)}

Test F1 recommendations in practice: - Select 10-15 high-confidence F1
pairs - Implement consolidation in pilot surveys - Measure data quality,
respondent burden, agency workflow

\textbf{Outcome}: Evidence on consolidation feasibility and impact

\subsection{Short-Term Extensions
(2026)}\label{short-term-extensions-2026}

\textbf{3. Expand Survey Coverage}

Apply methodology to additional surveys: - \textbf{SIPP} (Survey of
Income and Program Participation): 500+ questions - \textbf{NHIS}
(National Health Interview Survey): 400+ questions - \textbf{AHS}
(American Housing Survey): 300+ questions

\textbf{Outcome}: Comprehensive federal survey consolidation map

\textbf{4. Bidirectional Analysis}

Reverse the analysis: - Use CPS as consolidation target (test if FoodAPS
questions could adopt CPS measures) - Use FoodAPS as target (test
specialized question adoption)

\textbf{Outcome}: Identify opportunities for survey-specific measures to
propagate

\subsection{Long-Term Applications
(2027+)}\label{long-term-applications-2027}

\textbf{5. Survey Design Integration}

Incorporate harmonization analysis into questionnaire design: - Flag new
questions with consolidation potential - Recommend harmonized wording
for planned questions - Identify opportunities for question recycling

\textbf{Outcome}: Proactive burden reduction in survey development

\textbf{6. Cross-Survey Imputation}

Leverage consolidable questions for statistical integration: - Impute
FoodAPS responses using CPS data where F1/F2 relationships exist -
Enable cross-survey analysis without direct linkage - Validate
imputation accuracy using validation samples

\textbf{Outcome}: Enhanced data integration capabilities

\textbf{7. Respondent Burden Modeling}

Quantify burden reduction from consolidation: - Estimate cumulative
burden across surveys - Model impact of consolidating F1 questions -
Assess trade-offs between burden and data granularity

\textbf{Outcome}: Evidence-based burden reduction targets

\textbf{8. Methodology Refinement}

Incorporate expert feedback to improve: - Prompt engineering for higher
initial accuracy - Threshold optimization for triage quadrants - Barrier
taxonomy expansion for edge cases

\textbf{Outcome}: More accurate, efficient future analyses

\begin{center}\rule{0.5\linewidth}{0.5pt}\end{center}

\section{Broader Impact}\label{broader-impact}

\subsection{For the Federal Statistical
System}\label{for-the-federal-statistical-system}

This work provides: - \textbf{Data-driven foundation} for consolidation
decisions - \textbf{Reproducible methodology} applicable across surveys
- \textbf{Evidence} that AI can assist (not replace) expert judgment

\subsection{For Survey Methodology}\label{for-survey-methodology}

Demonstrates that: - Harmonization frameworks (DataSHaPER) apply beyond
epidemiology - Multi-model ensembles achieve high agreement on complex
classifications - Operational triage frameworks can reduce expert review
load

\subsection{For AI-Assisted Research}\label{for-ai-assisted-research}

Shows that: - LLMs can apply structured frameworks consistently -
Ensemble + arbitration patterns reduce single-model bias - Confidence
scores predict agreement and guide triage

\begin{center}\rule{0.5\linewidth}{0.5pt}\end{center}

\section{Closing Perspective}\label{closing-perspective}

Survey harmonization is fundamentally about \textbf{balancing
standardization with specialization}. Federal surveys serve distinct
research purposes, and consolidation should not come at the expense of
substantive content.

This analysis identifies where consolidation is \emph{technically
feasible} (44\% of questions). The decision to actually consolidate
requires: - \textbf{Stakeholder input}: Agency missions and research
priorities - \textbf{Data quality assessment}: Trade-offs between burden
and granularity - \textbf{Policy context}: Legislative mandates and user
needs

\textbf{AI assists, experts decide.} The methodology accelerates
analysis but does not replace judgment. By routing 75\% of questions to
auto-processing and 25\% to expert review, we focus human effort where
it matters most.

\begin{center}\rule{0.5\linewidth}{0.5pt}\end{center}

\section{Final Recommendations}\label{final-recommendations}

\subsection{For Agencies}\label{for-agencies}

\begin{enumerate}
\def\labelenumi{\arabic{enumi}.}
\tightlist
\item
  \textbf{Review F1 questions} (60 questions) - low-hanging fruit for
  consolidation
\item
  \textbf{Evaluate F2 questions} (108 questions) - assess cost-benefit
  of statistical adjustment
\item
  \textbf{Respect F3 questions} (212 questions) - maintain
  survey-specific content
\end{enumerate}

\subsection{For Methodologists}\label{for-methodologists}

\begin{enumerate}
\def\labelenumi{\arabic{enumi}.}
\tightlist
\item
  \textbf{Validate classifications} - expert review of 93 flagged
  questions
\item
  \textbf{Test consolidation pilots} - measure real-world impact
\item
  \textbf{Expand coverage} - apply methodology to additional surveys
\end{enumerate}

\subsection{For Data Users}\label{for-data-users-1}

\begin{enumerate}
\def\labelenumi{\arabic{enumi}.}
\tightlist
\item
  \textbf{Leverage harmonized variables} - enable cross-survey
  integration
\item
  \textbf{Assess fitness-for-purpose} - verify suitability for specific
  analyses
\item
  \textbf{Provide feedback} - inform consolidation priorities
\end{enumerate}

\begin{center}\rule{0.5\linewidth}{0.5pt}\end{center}

\section{Conclusion}\label{conclusion-1}

We set out to answer three questions:

\textbf{RQ1: What proportion of federal survey questions can be
consolidated?} - \textbf{Answer}: 44\% have at least one consolidable
ACS match (F1 or F2)

\textbf{RQ2: What barriers prevent consolidation?} - \textbf{Answer}:
97\% of failures stem from Construct/Concept differences, specifically
concept definitions (CC.1)

\textbf{RQ3: Can AI-assisted methods accelerate this analysis?} -
\textbf{Answer}: Yes - multi-model ensemble achieves κ=0.845 agreement
and reduces expert review load by 75\%

This work provides the federal statistical system with an evidence base
for consolidation decisions, a reproducible methodology for future
analyses, and specific recommendations for 380 questions. The findings
demonstrate that AI-assisted methods can accelerate traditionally
labor-intensive survey harmonization work while preserving expert
oversight and judgment.

\textbf{The path forward is clear}: Validate, pilot, expand. Survey
harmonization is no longer limited by analysis capacity - it's a matter
of prioritization and stakeholder alignment.

\begin{center}\rule{0.5\linewidth}{0.5pt}\end{center}

\textbf{Report complete. See Appendices for technical details, taxonomy
definitions, and expert review tables.}

\part{Appendices}

\chapter{Appendix A: Taxonomy
Definitions}\label{appendix-a-taxonomy-definitions}

\section{Harmonization Feasibility
Codes}\label{harmonization-feasibility-codes}

\subsection{F1: Direct Recode}\label{f1-direct-recode}

\textbf{Definition}: Variables are mechanically transformable through
simple recoding, collapsing categories, or unit conversion.

\textbf{Action Required}: Simple data transformation

\textbf{Characteristics}: - Same underlying construct - Compatible
measurement scale - Deterministic mapping between responses - No loss of
information (or acceptable minimal loss)

\textbf{Examples}:

\begin{longtable}[]{@{}
  >{\raggedright\arraybackslash}p{(\linewidth - 4\tabcolsep) * \real{0.2500}}
  >{\raggedright\arraybackslash}p{(\linewidth - 4\tabcolsep) * \real{0.2500}}
  >{\raggedright\arraybackslash}p{(\linewidth - 4\tabcolsep) * \real{0.5000}}@{}}
\toprule\noalign{}
\begin{minipage}[b]{\linewidth}\raggedright
Source
\end{minipage} & \begin{minipage}[b]{\linewidth}\raggedright
Target
\end{minipage} & \begin{minipage}[b]{\linewidth}\raggedright
Transformation
\end{minipage} \\
\midrule\noalign{}
\endhead
\bottomrule\noalign{}
\endlastfoot
Age in years (continuous) & Age in 5-year brackets & Group into
categories \\
Income in dollars & Income in thousands & Divide by 1,000 \\
Hours worked (0-168) & Hours worked (0-168) & Direct copy (identical) \\
Education (8 categories) & Education (4 categories) & Collapse
categories \\
\end{longtable}

\textbf{Data Quality}: F1 transformations preserve original information
or involve only aggregation (no modeling or imputation).

\begin{center}\rule{0.5\linewidth}{0.5pt}\end{center}

\subsection{F2: Statistical Adjustment}\label{f2-statistical-adjustment}

\textbf{Definition}: Variables require modeling, imputation, bridging
studies, or statistical harmonization methods to make comparable.

\textbf{Action Required}: Statistical harmonization (modeling,
imputation, or bridging)

\textbf{Characteristics}: - Related but not identical constructs -
Requires statistical assumptions - May involve measurement error or
information loss - Transformation is probabilistic, not deterministic

\textbf{Examples}:

\begin{longtable}[]{@{}
  >{\raggedright\arraybackslash}p{(\linewidth - 4\tabcolsep) * \real{0.2500}}
  >{\raggedright\arraybackslash}p{(\linewidth - 4\tabcolsep) * \real{0.2500}}
  >{\raggedright\arraybackslash}p{(\linewidth - 4\tabcolsep) * \real{0.5000}}@{}}
\toprule\noalign{}
\begin{minipage}[b]{\linewidth}\raggedright
Source
\end{minipage} & \begin{minipage}[b]{\linewidth}\raggedright
Target
\end{minipage} & \begin{minipage}[b]{\linewidth}\raggedright
Transformation
\end{minipage} \\
\midrule\noalign{}
\endhead
\bottomrule\noalign{}
\endlastfoot
Weekly work hours & Annual work weeks & Model seasonal patterns, impute
missing data \\
Hourly wage (excluding tips) & Hourly wage (including tips) & Impute tip
proportion based on occupation \\
Food insecurity (12-item battery) & Food insecurity (6-item battery) &
Bridging study to map scales \\
Health insurance (current) & Health insurance (annual) & Model coverage
transitions \\
\end{longtable}

\textbf{Data Quality}: F2 transformations introduce modeling assumptions
and potential bias. Users must assess fitness-for-purpose.

\begin{center}\rule{0.5\linewidth}{0.5pt}\end{center}

\subsection{F3: Incompatible}\label{f3-incompatible}

\textbf{Definition}: Variables measure fundamentally different
constructs and cannot be harmonized without re-fielding questions.

\textbf{Action Required}: No harmonization possible

\textbf{Characteristics}: - Different underlying constructs - No
statistical transformation can bridge gap - Conceptual rather than
operational barriers - Would require new data collection to resolve

\textbf{Examples}:

\begin{longtable}[]{@{}
  >{\raggedright\arraybackslash}p{(\linewidth - 4\tabcolsep) * \real{0.2353}}
  >{\raggedright\arraybackslash}p{(\linewidth - 4\tabcolsep) * \real{0.2353}}
  >{\raggedright\arraybackslash}p{(\linewidth - 4\tabcolsep) * \real{0.5294}}@{}}
\toprule\noalign{}
\begin{minipage}[b]{\linewidth}\raggedright
Source
\end{minipage} & \begin{minipage}[b]{\linewidth}\raggedright
Target
\end{minipage} & \begin{minipage}[b]{\linewidth}\raggedright
Why Incompatible
\end{minipage} \\
\midrule\noalign{}
\endhead
\bottomrule\noalign{}
\endlastfoot
``Do you want to work full-time?'' & ``Did you work last week?'' &
Intention vs.~behavior \\
``How many children live with you?'' & ``How many children have you
given birth to?'' & Household composition vs.~lifetime fertility \\
``Annual household income'' & ``Hourly wage at main job'' & Different
income concepts \\
``Food security score'' & ``SNAP participation'' & Outcome measure
vs.~program participation \\
\end{longtable}

\textbf{Barrier Codes}: F3 classifications must include a barrier code
explaining why harmonization is impossible.

\begin{center}\rule{0.5\linewidth}{0.5pt}\end{center}

\section{Barrier Code Taxonomy}\label{barrier-code-taxonomy}

When variables are F3 (incompatible), one of the following barrier codes
applies:

\subsection{Temporal Constraints (TC)}\label{temporal-constraints-tc-1}

Reference period or timing differences that prevent meaningful
comparison.

\subsubsection{TC.1: Reference Period
Length}\label{tc.1-reference-period-length}

\textbf{Definition}: Different recall windows or time spans.

\textbf{Examples}: - ``Hours worked last week'' vs.~``Hours worked in
past 12 months'' - ``Income last month'' vs.~``Income last year'' -
``Health status today'' vs.~``Health status past year''

\textbf{Why Incompatible}: Different reference periods capture different
phenomena (short-term vs.~long-term patterns, seasonal effects, recall
bias).

\begin{center}\rule{0.5\linewidth}{0.5pt}\end{center}

\subsubsection{TC.2: Temporal Framing}\label{tc.2-temporal-framing}

\textbf{Definition}: Point-in-time vs.~habitual vs.~retrospective
framing.

\textbf{Examples}: - ``Did you work last week?'' (point-in-time)
vs.~``Do you usually work?'' (habitual) - ``Are you currently
enrolled?'' (present) vs.~``Were you ever enrolled?'' (retrospective) -
``Last week's income'' vs.~``Usual weekly income''

\textbf{Why Incompatible}: Habitual measures smooth over variation;
point-in-time measures capture specific instances.

\begin{center}\rule{0.5\linewidth}{0.5pt}\end{center}

\subsubsection{TC.3: Calendar Alignment}\label{tc.3-calendar-alignment}

\textbf{Definition}: Fixed calendar periods vs.~rolling reference
periods.

\textbf{Examples}: - ``Income in calendar year 2025'' vs.~``Income in
past 12 months'' - ``Employment status January 1'' vs.~``Employment
status on interview date'' - ``Tax year'' vs.~``Fiscal year''

\textbf{Why Incompatible}: Calendar alignment affects comparability
across respondents interviewed at different times.

\begin{center}\rule{0.5\linewidth}{0.5pt}\end{center}

\subsection{Construct/Concept Constraints
(CC)}\label{constructconcept-constraints-cc}

Concept definition or operationalization differences - the most common
barrier type.

\subsubsection{CC.1: Concept Definition}\label{cc.1-concept-definition}

\textbf{Definition}: Different meaning of core term or different
construct entirely.

\textbf{Examples}: - ``Employment'' including vs.~excluding unpaid work
- ``Income'' including vs.~excluding government transfers -
``Household'' including vs.~excluding non-relatives - ``Disability''
based on self-report vs.~functional limitations

\textbf{Why Incompatible}: Measuring fundamentally different things,
even if using similar terminology.

\textbf{This is the dominant barrier} (70\% of F3 pairs).

\begin{center}\rule{0.5\linewidth}{0.5pt}\end{center}

\subsubsection{CC.2: Operationalization}\label{cc.2-operationalization}

\textbf{Definition}: Different behavioral indicators or measurement
approaches for nominally the same concept.

\textbf{Examples}: - ``Food insecurity'' measured by 12-item battery
vs.~single question - ``Health status'' measured by self-rating
vs.~diagnosed conditions - ``Job search'' measured by specific
activities vs.~general question - ``Poverty'' based on income thresholds
vs.~consumption measures

\textbf{Why Incompatible}: Different operationalizations capture
different dimensions of the construct.

\begin{center}\rule{0.5\linewidth}{0.5pt}\end{center}

\subsubsection{CC.3: Boundary
Conditions}\label{cc.3-boundary-conditions}

\textbf{Definition}: Different thresholds, cutoffs, or qualifying
conditions.

\textbf{Examples}: - ``Full-time work'' defined as 35+ hours vs.~40+
hours - ``Low income'' using different poverty thresholds - ``Recent
immigrant'' defined as \textless5 years vs.~\textless10 years -
``Elderly'' defined as 65+ vs.~60+

\textbf{Why Incompatible}: Boundary differences create non-overlapping
populations or classifications.

\begin{center}\rule{0.5\linewidth}{0.5pt}\end{center}

\subsubsection{CC.4: Scope Inclusions}\label{cc.4-scope-inclusions}

\textbf{Definition}: Different components counted or included in
measurement.

\textbf{Examples}: - ``Income'' including vs.~excluding capital gains -
``Household members'' including vs.~excluding temporary residents -
``Work hours'' including vs.~excluding overtime - ``Educational
attainment'' including vs.~excluding vocational training

\textbf{Why Incompatible}: Different scope creates systematically
different measures.

\begin{center}\rule{0.5\linewidth}{0.5pt}\end{center}

\subsection{Population/Coverage Constraints
(PC)}\label{populationcoverage-constraints-pc}

Universe, frame, or sample design differences.

\subsubsection{PC.1: Universe
Definition}\label{pc.1-universe-definition}

\textbf{Definition}: Target population differs between surveys.

\textbf{Examples}: - Civilian population vs.~all residents - Adults 18+
vs.~all household members - US citizens vs.~all residents - Employed
persons vs.~labor force participants

\textbf{Why Incompatible}: Questions asked of different populations
cannot be directly compared.

\begin{center}\rule{0.5\linewidth}{0.5pt}\end{center}

\subsubsection{PC.2: Frame Exclusions}\label{pc.2-frame-exclusions}

\textbf{Definition}: Different sampling frame exclusions.

\textbf{Examples}: - Excluding vs.~including group quarters (prisons,
dorms, nursing homes) - Excluding vs.~including US territories -
Excluding vs.~including institutionalized population - Excluding
vs.~including military bases

\textbf{Why Incompatible}: Systematic exclusions create non-comparable
samples.

\begin{center}\rule{0.5\linewidth}{0.5pt}\end{center}

\subsubsection{PC.3: Age Bounds}\label{pc.3-age-bounds}

\textbf{Definition}: Different age eligibility criteria.

\textbf{Examples}: - Labor force questions for 15+ vs.~16+ vs.~18+ -
Health insurance questions for all ages vs.~non-elderly only - Education
questions for school-age vs.~all adults

\textbf{Why Incompatible}: Age restrictions prevent full population
comparison.

\begin{center}\rule{0.5\linewidth}{0.5pt}\end{center}

\subsubsection{PC.4: Geographic Scope}\label{pc.4-geographic-scope}

\textbf{Definition}: Different geographic coverage.

\textbf{Examples}: - 50 states vs.~50 states + DC + territories -
Metropolitan areas only vs.~all areas - Border regions vs.~national
sample - State-specific vs.~national questions

\textbf{Why Incompatible}: Geographic restrictions limit comparability.

\begin{center}\rule{0.5\linewidth}{0.5pt}\end{center}

\subsection{Response Scale Constraints
(RS)}\label{response-scale-constraints-rs}

Scale type, categories, or format differences.

\subsubsection{RS.1: Scale Type}\label{rs.1-scale-type}

\textbf{Definition}: Fundamentally different response formats.

\textbf{Examples}: - Binary (yes/no) vs.~5-point Likert scale -
Continuous numeric vs.~categorical - Frequency scale vs.~intensity scale
- Open numeric vs.~bracketed ranges

\textbf{Why Incompatible}: Different scale types capture different
levels of information and are not directly comparable.

\begin{center}\rule{0.5\linewidth}{0.5pt}\end{center}

\subsubsection{RS.2: Category Structure}\label{rs.2-category-structure}

\textbf{Definition}: Different number of categories or category
boundaries within same scale type.

\textbf{Examples}: - 5-point vs.~7-point Likert scale - Income brackets
with different boundaries - Age categories with different groupings -
``Often/Sometimes/Rarely'' vs.~``Always/Often/Sometimes/Rarely/Never''

\textbf{Why Incompatible}: Category differences prevent direct mapping
without information loss or assumptions.

\begin{center}\rule{0.5\linewidth}{0.5pt}\end{center}

\subsubsection{RS.3: Anchoring/Labels}\label{rs.3-anchoringlabels}

\textbf{Definition}: Different verbal anchors or category labels.

\textbf{Examples}: - ``Strongly agree'' interpreted differently across
surveys - ``Good health'' vs.~``Excellent health'' as top category -
``Frequently'' without numeric anchors - Cultural/linguistic differences
in label interpretation

\textbf{Why Incompatible}: Verbal anchors are subjective and
survey-specific.

\begin{center}\rule{0.5\linewidth}{0.5pt}\end{center}

\subsubsection{RS.4: Numeric vs.~Verbal}\label{rs.4-numeric-vs.-verbal}

\textbf{Definition}: Numeric scales vs.~labeled categories.

\textbf{Examples}: - 0-10 slider vs.~``Excellent/Very
Good/Good/Fair/Poor'' - Days per week (numeric)
vs.~``Daily/Weekly/Monthly'' (verbal) - Exact count vs.~frequency
categories

\textbf{Why Incompatible}: Numeric precision differs from categorical
judgment.

\begin{center}\rule{0.5\linewidth}{0.5pt}\end{center}

\subsection{Mode/Context Constraints
(MC)}\label{modecontext-constraints-mc}

Interview mode or questionnaire context differences.

\subsubsection{MC.1: Interview Mode}\label{mc.1-interview-mode}

\textbf{Definition}: Different data collection modes.

\textbf{Examples}: - CATI (phone) vs.~Web self-administered - In-person
interview vs.~mail survey - Interviewer-administered
vs.~self-administered - Video vs.~audio interview

\textbf{Why Incompatible}: Mode effects (social desirability,
interviewer effects, visual cues) create systematic differences.

\begin{center}\rule{0.5\linewidth}{0.5pt}\end{center}

\subsubsection{MC.2: Question Routing}\label{mc.2-question-routing}

\textbf{Definition}: Different skip patterns or conditional questions.

\textbf{Examples}: - Asked of all respondents vs.~subset - Follow-up
question vs.~standalone question - Conditional on prior response
vs.~unconditional

\textbf{Why Incompatible}: Routing affects who answers and question
interpretation context.

\begin{center}\rule{0.5\linewidth}{0.5pt}\end{center}

\subsubsection{MC.3: Contextual Priming}\label{mc.3-contextual-priming}

\textbf{Definition}: Preceding questions affect interpretation.

\textbf{Examples}: - Question order effects - Section context (health
questions in health module vs.~employment module) - Priming from related
questions - Survey frame (government survey vs.~academic survey)

\textbf{Why Incompatible}: Context shapes interpretation and response
patterns.

\begin{center}\rule{0.5\linewidth}{0.5pt}\end{center}

\subsubsection{MC.4: Proxy Response}\label{mc.4-proxy-response}

\textbf{Definition}: Proxy vs.~self-report rules differ.

\textbf{Examples}: - Self-report only vs.~household head can report for
all - Proxy allowed with restrictions vs.~always allowed - Age threshold
for self-report differs

\textbf{Why Incompatible}: Proxy responses differ systematically from
self-reports.

\begin{center}\rule{0.5\linewidth}{0.5pt}\end{center}

\subsection{Processing/Metadata Constraints
(PM)}\label{processingmetadata-constraints-pm}

Coding, weighting, or documentation differences.

\subsubsection{PM.1: Coding Schemes}\label{pm.1-coding-schemes}

\textbf{Definition}: Different classification systems.

\textbf{Examples}: - Occupation coding differences (SOC 2010 vs.~SOC
2018) - Industry coding (NAICS 2012 vs.~NAICS 2017) - Geographic coding
(FIPS vs.~other systems) - Race/ethnicity categories (different
standards)

\textbf{Why Incompatible}: Different coding prevents direct comparison
without crosswalks.

\begin{center}\rule{0.5\linewidth}{0.5pt}\end{center}

\subsubsection{PM.2: Derived Variables}\label{pm.2-derived-variables}

\textbf{Definition}: Different construction algorithms.

\textbf{Examples}: - Poverty calculation methods (official
vs.~supplemental) - Composite scores with different formulas - Imputed
variables with different models - Weighted vs.~unweighted aggregates

\textbf{Why Incompatible}: Different construction creates systematically
different measures.

\begin{center}\rule{0.5\linewidth}{0.5pt}\end{center}

\subsubsection{PM.3: Documentation Gaps}\label{pm.3-documentation-gaps}

\textbf{Definition}: Insufficient metadata to assess comparability.

\textbf{Examples}: - Missing questionnaire context - Unclear skip
patterns - Ambiguous variable definitions - Undocumented changes across
survey years

\textbf{Why Incompatible}: Cannot determine compatibility without
adequate documentation.

\begin{center}\rule{0.5\linewidth}{0.5pt}\end{center}

\section{Application Notes}\label{application-notes}

\subsection{Selecting Primary Barrier
Code}\label{selecting-primary-barrier-code}

When multiple barriers apply, select based on hierarchy: 1. \textbf{CC}
(Construct) - most fundamental 2. \textbf{TC} (Temporal) - affects
interpretation 3. \textbf{RS} (Response Scale) - limits comparability 4.
\textbf{PC} (Population) - affects universe 5. \textbf{MC}
(Mode/Context) - measurement effects 6. \textbf{PM} (Processing) -
technical issues

\textbf{Example}: If a question has both CC.1 (different concept) and
RS.1 (different scale), classify as CC.1 because construct differences
are more fundamental.

\subsection{Edge Cases}\label{edge-cases}

\textbf{Borderline F2/F3}: When uncertain whether statistical adjustment
can bridge gap: - \textbf{F2} if adjustment is feasible with reasonable
assumptions - \textbf{F3} if adjustment requires untenable assumptions
or would destroy information

\textbf{Multiple sub-codes}: Document all applicable codes in reasoning
field, but select single primary code for classification.

\begin{center}\rule{0.5\linewidth}{0.5pt}\end{center}

\section{Citation}\label{citation}

This taxonomy is adapted from:

\begin{itemize}
\item
  \textbf{Fortier, I., et al.~(2011)}. Is rigorous retrospective
  harmonization possible? Application of the DataSHaPER approach across
  53 large studies. \emph{International Journal of Epidemiology}, 40(5),
  1314-1328.
\item
  \textbf{Fortier, I., et al.~(2017)}. Maelstrom Research guidelines for
  rigorous retrospective data harmonization. \emph{International Journal
  of Epidemiology}, 46(1), 103-115.
\end{itemize}

Extensions and examples specific to federal surveys developed for this
analysis.

\chapter{Appendix B: Methodology
Decisions}\label{appendix-b-methodology-decisions}

This appendix documents key methodological decisions made during the
analysis, including rationale and alternatives considered.

\begin{center}\rule{0.5\linewidth}{0.5pt}\end{center}

\section{Decision 001: Three-Model
Ensemble}\label{decision-001-three-model-ensemble}

\textbf{Date}: 2026-01-28

\textbf{Decision}: Use three frontier LLMs (OpenAI gpt-4o-mini,
Anthropic claude-3-5-haiku, Google gemini-2.0-flash-exp) as independent
raters.

\textbf{Rationale}: - Reduces single-model bias - Different training
data and architectures provide diversity - Ensemble agreement validates
consistency of judgments

\textbf{Alternatives Considered}: - Single model only - rejected due to
bias risk - Five+ models - rejected due to cost and diminishing returns

\textbf{Outcome}: High inter-rater agreement (κ = 0.845) validated
approach.

\begin{center}\rule{0.5\linewidth}{0.5pt}\end{center}

\section{Decision 002: Structured JSON
Output}\label{decision-002-structured-json-output}

\textbf{Date}: 2026-01-28

\textbf{Decision}: Require models to output structured JSON with
feasibility, barrier\_code, confidence, and reasoning fields.

\textbf{Rationale}: - Enables automated parsing and aggregation - Forces
models to provide explicit reasoning - Confidence field supports triage

\textbf{Alternatives Considered}: - Free-text output - rejected due to
parsing complexity - Single classification only - rejected as
insufficient for expert review

\textbf{Outcome}: 100\% successful parsing with robust extraction
strategies.

\begin{center}\rule{0.5\linewidth}{0.5pt}\end{center}

\section{Decision 003: Parallel Processing with
Checkpointing}\label{decision-003-parallel-processing-with-checkpointing}

\textbf{Date}: 2026-01-28

\textbf{Decision}: Process pairs in parallel (6 workers) with
checkpoint/resume capability.

\textbf{Rationale}: - Reduces runtime (hours vs.~days) - Allows graceful
recovery from API failures - Supports cost control (stop/resume as
needed)

\textbf{Implementation}: - Batch size: 10 pairs per API call -
Exponential backoff for rate limits - JSON checkpoint file tracks
progress

\textbf{Outcome}: Robust execution with zero data loss.

\begin{center}\rule{0.5\linewidth}{0.5pt}\end{center}

\section{Decision 004: Google Rate Limit
Acceptance}\label{decision-004-google-rate-limit-acceptance}

\textbf{Date}: 2026-01-29

\textbf{Decision}: Proceed with OpenAI-Anthropic two-way comparison
after Google rate-limited at 751 pairs (47\%).

\textbf{Rationale}: - Two-way agreement sufficient for validation (κ =
0.845) - Google subset (751 pairs) confirms consistency (κ = 0.833) -
Cost/benefit of waiting for Google rate limit unfavorable

\textbf{Alternatives Considered}: - Wait for rate limit reset - rejected
due to timeline - Replace Google with another model - rejected as Google
data still valuable validation

\textbf{Outcome}: Primary findings based on OpenAI-Anthropic, Google
used for validation subset.

\begin{center}\rule{0.5\linewidth}{0.5pt}\end{center}

\section{Decision 005: Arbitration for Disagreements
Only}\label{decision-005-arbitration-for-disagreements-only}

\textbf{Date}: 2026-01-29

\textbf{Decision}: Arbitrate only disagreements (339 pairs, 21.2\%), not
all pairs.

\textbf{Rationale}: - Cost efficiency (arbitration uses stronger models)
- High-agreement pairs don't benefit from arbitration - Focuses
arbitrator attention on genuinely uncertain cases

\textbf{Alternatives Considered}: - Arbitrate all pairs - rejected due
to cost - No arbitration - rejected as disagreements need resolution

\textbf{Outcome}: Arbitration resolved 339 pairs with 87.6\% arbitrator
agreement.

\begin{center}\rule{0.5\linewidth}{0.5pt}\end{center}

\section{Decision 006: Question-Level Rollup (Not Pair-Level
Reporting)}\label{decision-006-question-level-rollup-not-pair-level-reporting}

\textbf{Date}: 2026-01-30

\textbf{Decision}: Report primary findings at question level (380
questions), not pair level (1,598 pairs).

\textbf{Rationale}: - Pair-level inflates denominator (multiple
comparisons per question) - Question-level answers stakeholder question:
``Is this question consolidable?'' - Statistical dependencies at pair
level violate independence assumptions

\textbf{Alternatives Considered}: - Pair-level only - rejected as
misleading (one F1 match makes question consolidable regardless of other
pairs) - Report both - \textbf{SELECTED}: Pair-level for research
detail, question-level for stakeholder summary

\textbf{Outcome}: Primary findings at question level (44.2\%
consolidable), pair-level data in appendix.

\begin{center}\rule{0.5\linewidth}{0.5pt}\end{center}

\section{Decision 007: Best-Match Selection per
Question}\label{decision-007-best-match-selection-per-question}

\textbf{Date}: 2026-01-30

\textbf{Decision}: For each source question, select single best ACS
match using hierarchy: F1 \textgreater{} F2 \textgreater{} F3, then
highest Borda score within tie.

\textbf{Rationale}: - Provides actionable recommendation (one pairing
per question) - Reflects real-world use case (each question maps to at
most one alternative) - Simplifies expert review table

\textbf{Alternatives Considered}: - Show all matches - rejected as
overwhelming for stakeholders - Average across matches - rejected as
statistically inappropriate (dependencies)

\textbf{Outcome}: 380 question-level records with clear best-match
recommendation.

\begin{center}\rule{0.5\linewidth}{0.5pt}\end{center}

\section{Decision 008: Four Scoring Methods (Ensemble
Approach)}\label{decision-008-four-scoring-methods-ensemble-approach}

\textbf{Date}: 2026-01-30

\textbf{Decision}: Compute four complementary scores: Composite,
Entropy, Bayesian, Borda.

\textbf{Rationale}: - Different scores capture different dimensions: -
Composite: Baseline (unanimous/majority) - Entropy: Stability (how much
they argued) - Bayesian: Probabilistic (with prior beliefs) - Borda:
Direction (leaning consolidable or not) - Ensemble avoids over-reliance
on single metric

\textbf{Alternatives Considered}: - Single score only - rejected as
insufficient to capture complexity - More than four - rejected as
diminishing returns

\textbf{Outcome}: Borda and Entropy selected for triage (Decision 016).

\begin{center}\rule{0.5\linewidth}{0.5pt}\end{center}

\section{Decision 009: Median Split for
Thresholds}\label{decision-009-median-split-for-thresholds}

\textbf{Date}: 2026-01-30

\textbf{Decision}: Use median split on question-level best-match scores
to define Borda and Entropy thresholds.

\textbf{Rationale}: - Pragmatic, interpretable (50th percentile) -
Avoids arbitrary threshold selection - Creates balanced quadrants for
triage

\textbf{Values}: - Borda median: 0.167 - Entropy median: 0.330

\textbf{Alternatives Considered}: - Optimize thresholds - rejected as no
ground truth for optimization - Fixed thresholds (0.5, 0.75) - rejected
as not data-driven

\textbf{Outcome}: Balanced quadrant distribution (Q1=151, Q2=136, Q3=40,
Q4=53).

\begin{center}\rule{0.5\linewidth}{0.5pt}\end{center}

\section{Decision 010: Question-Level Medians (Not
Pair-Level)}\label{decision-010-question-level-medians-not-pair-level}

\textbf{Date}: 2026-01-30

\textbf{Decision}: Compute threshold medians from 380 question-level
scores, not 1,598 pair-level scores.

\textbf{Rationale}: - Pair-level distribution dominated by unanimous F3
pairs (median Borda ≈ 0) - Question-level reflects actual triage
decision space - Threshold goal is to separate questions, not pairs

\textbf{Alternatives Considered}: - Pair-level medians - rejected as
produces degenerate thresholds (Borda median ≈ 0) - Fixed thresholds -
rejected as not data-driven

\textbf{Outcome}: See \texttt{docs/stage4\_ensemble\_methodology.md}
``Threshold Computation'' section for detailed justification.

\begin{center}\rule{0.5\linewidth}{0.5pt}\end{center}

\section{Decision 016: Two-Axis Triage Framework (Borda ×
Entropy)}\label{decision-016-two-axis-triage-framework-borda-entropy}

\textbf{Date}: 2026-01-31

\textbf{Decision}: Use Borda (direction) and Entropy (stability) as
two-axis triage framework for expert review routing.

\textbf{Rationale}: - \textbf{Borda}: Captures ensemble direction
(consolidable or not) - \textbf{Entropy}: Captures ensemble stability
(agreement or disagreement) - \textbf{Orthogonal dimensions}: Separates
``what's the answer?'' from ``how certain are we?''

\textbf{Quadrant Interpretation}: - Q1 (High Borda, High Entropy):
Confident consolidable → auto-accept - Q2 (Low Borda, High Entropy):
Confident non-consolidable → auto-reject - Q3 (High Borda, Low Entropy):
Uncertain accept (leaning yes but unstable) → \textbf{expert priority} -
Q4 (Low Borda, Low Entropy): Uncertain reject (ambiguous) → expert
secondary

\textbf{Alternatives Considered}: 1. Composite score only - rejected as
single dimension insufficient 2. Entropy only - rejected as doesn't
indicate direction 3. Bayesian + Borda - rejected as Bayesian highly
correlated with Borda (r=0.95) 4. \textbf{Borda + Entropy (SELECTED)} -
orthogonal (r=0.26), complementary dimensions

\textbf{Correlation Analysis}:

\begin{verbatim}
          Composite  Entropy  Bayesian  Borda
Composite    1.000    0.398     0.955   0.945
Entropy      0.398    1.000     0.341   0.260
Bayesian     0.955    0.341     1.000   0.950
Borda        0.945    0.260     0.950   1.000
\end{verbatim}

\textbf{Key Insight}: Borda and Entropy are least correlated (r=0.26),
maximizing information from two-axis framework.

\textbf{Outcome}: - 76\% auto-processed (Q1 + Q2) - 24\% expert review
(Q3 + Q4) - Q3 prioritized for review (uncertain accept cases)

\textbf{Framing}: This is an \textbf{operational tool}, not a
theoretical contribution. We do not claim novelty - the approach is a
pragmatic application of ensemble scoring for triage purposes. See
\texttt{docs/stage4\_ensemble\_methodology.md} for extended discussion
with ``sober framing''.

\textbf{Documentation}: Full decision rationale, correlation analysis,
and literature review in
\texttt{docs/methodology\_log\_decision\_016.md}.

\begin{center}\rule{0.5\linewidth}{0.5pt}\end{center}

\section{Decision 017: Sober Framing for Two-Axis
Approach}\label{decision-017-sober-framing-for-two-axis-approach}

\textbf{Date}: 2026-01-31

\textbf{Decision}: Frame two-axis triage as operational tool, not
theoretical contribution.

\textbf{Rationale}: - Approach is pragmatic heuristic for this project -
No claim of optimality or theoretical novelty - Avoid over-claiming in
research framing

\textbf{Language Used}: - ``Operational framework'' not ``novel
methodology'' - ``Useful heuristic'' not ``discovery'' - ``Performed
well in this context'' not ``best practice''

\textbf{Contrast}: - \textbf{Overstated}: ``We introduce a novel
two-axis framework\ldots{}'' - \textbf{Appropriate}: ``We use a two-axis
triage approach (Borda direction × Entropy stability) as an operational
tool\ldots{}''

\textbf{Documentation}: - Technical details in methodology section -
``Future Exploration'' section in
\texttt{stage4\_ensemble\_methodology.md} for research threads -
Citation queries in \texttt{docs/citation\_queries\_decision\_016.md}
remain available if needed

\textbf{Outcome}: Clear, honest framing that serves operational needs
without inflated claims.

\begin{center}\rule{0.5\linewidth}{0.5pt}\end{center}

\section{Decision 018: Expert Review Tables as Primary
Deliverable}\label{decision-018-expert-review-tables-as-primary-deliverable}

\textbf{Date}: 2026-01-31

\textbf{Decision}: Generate CSV tables (expert\_review\_cps.csv,
expert\_review\_foodaps.csv, expert\_review\_combined.csv) as primary
stakeholder deliverable.

\textbf{Rationale}: - Stakeholders need actionable, reviewable data -
CSV format enables filtering, sorting, annotation - Includes all context
(question text, match text, scores, reasoning)

\textbf{Columns Included}: - Source question ID and text - Best ACS
match ID and text - Feasibility classification (F1/F2/F3) - Barrier code
(if F3) - Scores (Borda, Entropy) - Triage quadrant (Q1-Q4) - Reasoning
summary

\textbf{Outcome}: 380 rows ready for expert review and validation.

\begin{center}\rule{0.5\linewidth}{0.5pt}\end{center}

\section{Decision 019: Visualization
Strategy}\label{decision-019-visualization-strategy}

\textbf{Date}: 2026-02-02

\textbf{Decision}: Generate both pair-level and question-level
visualizations.

\textbf{Rationale}: - \textbf{Pair-level} (1,598 pairs): Shows research
rigor (exhaustive comparison) - \textbf{Question-level} (380 questions):
Shows practical outcome (consolidation candidates) - Both perspectives
validate each other and serve different audiences

\textbf{Visualizations Created}: 1. Consolidation rates by survey
(question-level) 2. Barrier distribution (pair-level) 3. Expert review
load (question-level) 4. Harmonization code distribution (pair-level) 5.
Question consolidation distribution (question-level by survey)

\textbf{Outcome}: Comprehensive visual evidence for both research and
stakeholder audiences.

\begin{center}\rule{0.5\linewidth}{0.5pt}\end{center}

\section{Decision 020: Pipeline
Integration}\label{decision-020-pipeline-integration}

\textbf{Date}: 2026-01-31

\textbf{Decision}: Integrate all stages into \texttt{run\_pipeline.py}
orchestrator with resume capability.

\textbf{Rationale}: - Reproducibility requires unified execution -
Resume capability handles failures gracefully - Documentation through
code (pipeline is the specification)

\textbf{Implementation}: - Stage dependencies checked automatically -
Output validation before proceeding - Command-line flags:
\texttt{-\/-clean}, \texttt{-\/-from}, \texttt{-\/-only}

\textbf{Outcome}: Full pipeline runnable via single command:
\texttt{python\ run\_pipeline.py}

\begin{center}\rule{0.5\linewidth}{0.5pt}\end{center}

\section{Summary of Key Principles}\label{summary-of-key-principles}

Across all decisions, several principles guided methodology:

\begin{enumerate}
\def\labelenumi{\arabic{enumi}.}
\tightlist
\item
  \textbf{Transparency}: Document rationale, alternatives, and
  trade-offs
\item
  \textbf{Reproducibility}: Version control, checkpointing,
  deterministic execution
\item
  \textbf{Validation}: Inter-rater agreement, quality checks, expert
  review
\item
  \textbf{Pragmatism}: Operational tools over theoretical claims
\item
  \textbf{Stakeholder Focus}: Deliverables answer real questions (``Is
  this question consolidable?'')
\end{enumerate}

These principles ensured that analysis serves its intended purpose:
\textbf{providing evidence for federal survey consolidation decisions}.

\begin{center}\rule{0.5\linewidth}{0.5pt}\end{center}

\textbf{For complete decision log, see}:
\texttt{docs/methodology\_log.md} (all decisions) and
\texttt{docs/methodology\_log\_decision\_016.md} (detailed Decision 016
documentation).

\chapter{Appendix C: Expert Review
Tables}\label{appendix-c-expert-review-tables}

This appendix describes the expert review deliverables generated from
the analysis.

\begin{center}\rule{0.5\linewidth}{0.5pt}\end{center}

\section{Overview}\label{overview-1}

Three CSV files provide complete question-level results for expert
review and validation:

\begin{longtable}[]{@{}
  >{\raggedright\arraybackslash}p{(\linewidth - 4\tabcolsep) * \real{0.2400}}
  >{\raggedright\arraybackslash}p{(\linewidth - 4\tabcolsep) * \real{0.2400}}
  >{\raggedright\arraybackslash}p{(\linewidth - 4\tabcolsep) * \real{0.5200}}@{}}
\toprule\noalign{}
\begin{minipage}[b]{\linewidth}\raggedright
File
\end{minipage} & \begin{minipage}[b]{\linewidth}\raggedright
Rows
\end{minipage} & \begin{minipage}[b]{\linewidth}\raggedright
Description
\end{minipage} \\
\midrule\noalign{}
\endhead
\bottomrule\noalign{}
\endlastfoot
\texttt{expert\_review\_cps.csv} & 240 & CPS questions with ACS
matches \\
\texttt{expert\_review\_foodaps.csv} & 140 & FoodAPS questions with ACS
matches \\
\texttt{expert\_review\_combined.csv} & 380 & All questions combined \\
\end{longtable}

\textbf{Location}:
\texttt{reports/03\_harmonization\_constraints/output/analysis/}

\begin{center}\rule{0.5\linewidth}{0.5pt}\end{center}

\section{Table Schema}\label{table-schema}

Each file contains 17 columns:

\subsection{Identification}\label{identification}

\begin{longtable}[]{@{}lll@{}}
\toprule\noalign{}
Column & Type & Description \\
\midrule\noalign{}
\endhead
\bottomrule\noalign{}
\endlastfoot
\texttt{survey} & string & Source survey (CPS or FOODAPS) \\
\texttt{source\_q\_id} & string & Source question identifier \\
\texttt{source\_text} & string & Full source question text \\
\texttt{best\_match\_q\_id} & string & Best matching ACS question
identifier \\
\texttt{best\_match\_text} & string & Full ACS question text \\
\texttt{pair\_id} & string & Unique pair identifier (source\_target) \\
\end{longtable}

\subsection{Classification}\label{classification}

\begin{longtable}[]{@{}lll@{}}
\toprule\noalign{}
Column & Type & Description \\
\midrule\noalign{}
\endhead
\bottomrule\noalign{}
\endlastfoot
\texttt{best\_feasibility} & string & F1, F2, or F3 \\
\texttt{has\_consolidable\_path} & boolean & True if F1 or F2 \\
\texttt{barrier\_code} & string & CC.1, TC.2, etc. (if F3) \\
\texttt{confidence} & string & HIGH, MODERATE, or LOW \\
\end{longtable}

\subsection{Scoring}\label{scoring}

\begin{longtable}[]{@{}lll@{}}
\toprule\noalign{}
Column & Type & Description \\
\midrule\noalign{}
\endhead
\bottomrule\noalign{}
\endlastfoot
\texttt{score\_borda} & float & Borda count {[}0, 1{]} - direction \\
\texttt{score\_entropy} & float & Normalized entropy {[}0, 1{]} -
stability \\
\texttt{score\_composite} & float & Baseline composite score \\
\texttt{score\_bayesian} & float & Bayesian score with priors \\
\end{longtable}

\subsection{Triage}\label{triage}

\begin{longtable}[]{@{}lll@{}}
\toprule\noalign{}
Column & Type & Description \\
\midrule\noalign{}
\endhead
\bottomrule\noalign{}
\endlastfoot
\texttt{triage\_quadrant} & string & Q1, Q2, Q3, or Q4 \\
\end{longtable}

\subsection{Reasoning}\label{reasoning}

\begin{longtable}[]{@{}
  >{\raggedright\arraybackslash}p{(\linewidth - 4\tabcolsep) * \real{0.2963}}
  >{\raggedright\arraybackslash}p{(\linewidth - 4\tabcolsep) * \real{0.2222}}
  >{\raggedright\arraybackslash}p{(\linewidth - 4\tabcolsep) * \real{0.4815}}@{}}
\toprule\noalign{}
\begin{minipage}[b]{\linewidth}\raggedright
Column
\end{minipage} & \begin{minipage}[b]{\linewidth}\raggedright
Type
\end{minipage} & \begin{minipage}[b]{\linewidth}\raggedright
Description
\end{minipage} \\
\midrule\noalign{}
\endhead
\bottomrule\noalign{}
\endlastfoot
\texttt{reasoning\_summary} & string & Explanation of classification \\
\texttt{arbitration\_notes} & string & Additional notes from arbitration
(if applicable) \\
\end{longtable}

\begin{center}\rule{0.5\linewidth}{0.5pt}\end{center}

\section{Triage Quadrants}\label{triage-quadrants}

Questions are assigned to quadrants for prioritization:

\subsection{Q1: Confident Consolidable (151
questions)}\label{q1-confident-consolidable-151-questions}

\begin{itemize}
\tightlist
\item
  \textbf{Borda}: High (≥0.167)
\item
  \textbf{Entropy}: High (≥0.330)
\item
  \textbf{Interpretation}: Models agree it's consolidable
\item
  \textbf{Action}: Auto-accept, spot-check recommended
\item
  \textbf{Expert Review Priority}: Low
\end{itemize}

\textbf{Example}:

\begin{verbatim}
Source (CPS): "How many hours per week do you USUALLY work?"
Target (ACS): "What is your best estimate of hours per week you usually work?"
Feasibility: F1
Borda: 1.0, Entropy: 1.0
\end{verbatim}

\begin{center}\rule{0.5\linewidth}{0.5pt}\end{center}

\subsection{Q2: Confident Non-Consolidable (136
questions)}\label{q2-confident-non-consolidable-136-questions}

\begin{itemize}
\tightlist
\item
  \textbf{Borda}: Low (\textless0.167)
\item
  \textbf{Entropy}: High (≥0.330)
\item
  \textbf{Interpretation}: Models agree it's not consolidable
\item
  \textbf{Action}: Auto-reject, low priority for review
\item
  \textbf{Expert Review Priority}: Low
\end{itemize}

\textbf{Example}:

\begin{verbatim}
Source (CPS): "Do you WANT to work full-time (35+ hours)?"
Target (ACS): "Did you work for pay last week?"
Feasibility: F3 (CC.1)
Borda: 0.0, Entropy: 0.33
\end{verbatim}

\begin{center}\rule{0.5\linewidth}{0.5pt}\end{center}

\subsection{Q3: Uncertain Accept (40 questions) ⚠️
PRIORITY}\label{q3-uncertain-accept-40-questions-priority}

\begin{itemize}
\tightlist
\item
  \textbf{Borda}: High (≥0.167)
\item
  \textbf{Entropy}: Low (\textless0.330)
\item
  \textbf{Interpretation}: Leaning consolidable but unstable (models
  argued)
\item
  \textbf{Action}: \textbf{EXPERT REVIEW REQUIRED} - highest priority
\item
  \textbf{Expert Review Priority}: \textbf{HIGH}
\end{itemize}

\textbf{Why Priority?} - High Borda suggests consolidation potential -
Low Entropy indicates disagreement or uncertainty - These are ``edge
cases'' that need human judgment

\textbf{Review Question}: Is consolidation truly feasible despite model
disagreement?

\begin{center}\rule{0.5\linewidth}{0.5pt}\end{center}

\subsection{Q4: Uncertain Reject (53
questions)}\label{q4-uncertain-reject-53-questions}

\begin{itemize}
\tightlist
\item
  \textbf{Borda}: Low (\textless0.167)
\item
  \textbf{Entropy}: Low (\textless0.330)
\item
  \textbf{Interpretation}: Genuinely ambiguous (no clear direction, low
  agreement)
\item
  \textbf{Action}: Expert review - secondary priority
\item
  \textbf{Expert Review Priority}: Medium
\end{itemize}

\textbf{Why Secondary?} - Low Borda suggests non-consolidable lean - Low
Entropy indicates genuine ambiguity - Less actionable than Q3 (less
clear consolidation potential)

\textbf{Review Question}: Can expert domain knowledge resolve ambiguity?

\begin{center}\rule{0.5\linewidth}{0.5pt}\end{center}

\section{Usage Guide}\label{usage-guide}

\subsection{For Expert Reviewers}\label{for-expert-reviewers}

\textbf{Step 1: Prioritize Q3 Questions (40 questions)}

Focus first on Q3 - these have consolidation potential but model
uncertainty:

\begin{Shaded}
\begin{Highlighting}[]
\CommentTok{\# Filter Q3 questions}
\FunctionTok{grep} \StringTok{",Q3,"}\NormalTok{ expert\_review\_combined.csv }\OperatorTok{\textgreater{}}\NormalTok{ q3\_priority\_review.csv}
\end{Highlighting}
\end{Shaded}

\textbf{Review Criteria}: - Does source question measure same construct
as target? - Are differences fixable (F2) or fundamental (F3)? - Is
consolidation desirable from policy/research perspective?

\textbf{Actions}: - \textbf{Confirm F1/F2}: Validate consolidation is
appropriate - \textbf{Reclassify to F3}: If models missed fundamental
barrier - \textbf{Add notes}: Document reasoning for stakeholders

\begin{center}\rule{0.5\linewidth}{0.5pt}\end{center}

\textbf{Step 2: Review Q4 Questions (53 questions)}

After Q3, review Q4 - genuinely ambiguous cases:

\begin{Shaded}
\begin{Highlighting}[]
\CommentTok{\# Filter Q4 questions}
\FunctionTok{grep} \StringTok{",Q4,"}\NormalTok{ expert\_review\_combined.csv }\OperatorTok{\textgreater{}}\NormalTok{ q4\_secondary\_review.csv}
\end{Highlighting}
\end{Shaded}

\textbf{Review Criteria}: - Why did models disagree? - Is additional
context needed? - Does domain expertise resolve ambiguity?

\begin{center}\rule{0.5\linewidth}{0.5pt}\end{center}

\textbf{Step 3: Spot-Check Q1/Q2 (287 questions)}

Sample Q1 and Q2 to validate auto-processing quality:

\begin{Shaded}
\begin{Highlighting}[]
\CommentTok{\# Sample 10\% of Q1 questions}
\FunctionTok{grep} \StringTok{",Q1,"}\NormalTok{ expert\_review\_combined.csv }\KeywordTok{|} \FunctionTok{shuf} \AttributeTok{{-}n}\NormalTok{ 15 }\OperatorTok{\textgreater{}}\NormalTok{ q1\_spot\_check.csv}

\CommentTok{\# Sample 10\% of Q2 questions}
\FunctionTok{grep} \StringTok{",Q2,"}\NormalTok{ expert\_review\_combined.csv }\KeywordTok{|} \FunctionTok{shuf} \AttributeTok{{-}n}\NormalTok{ 14 }\OperatorTok{\textgreater{}}\NormalTok{ q2\_spot\_check.csv}
\end{Highlighting}
\end{Shaded}

\textbf{Review Criteria}: - Are high-confidence consolidable (Q1)
classifications correct? - Are high-confidence non-consolidable (Q2)
classifications correct? - Any systematic errors that suggest prompt
refinement?

\begin{center}\rule{0.5\linewidth}{0.5pt}\end{center}

\subsection{For Data Users}\label{for-data-users-2}

\textbf{Finding Consolidable Questions}:

\begin{Shaded}
\begin{Highlighting}[]
\CommentTok{\# All consolidable questions (F1 + F2)}
\FunctionTok{grep} \StringTok{",True,"}\NormalTok{ expert\_review\_combined.csv }\OperatorTok{\textgreater{}}\NormalTok{ consolidable\_questions.csv}

\CommentTok{\# Direct recode only (F1)}
\FunctionTok{grep} \StringTok{",F1,"}\NormalTok{ expert\_review\_combined.csv }\OperatorTok{\textgreater{}}\NormalTok{ f1\_direct\_recode.csv}

\CommentTok{\# Statistical adjustment (F2)}
\FunctionTok{grep} \StringTok{",F2,"}\NormalTok{ expert\_review\_combined.csv }\OperatorTok{\textgreater{}}\NormalTok{ f2\_statistical\_adjustment.csv}
\end{Highlighting}
\end{Shaded}

\textbf{Finding Specific Topics}:

\begin{Shaded}
\begin{Highlighting}[]
\CommentTok{\# Employment questions (example)}
\FunctionTok{grep} \AttributeTok{{-}i} \StringTok{"work\textbackslash{}|employ\textbackslash{}|job\textbackslash{}|hour"}\NormalTok{ expert\_review\_combined.csv }\OperatorTok{\textgreater{}}\NormalTok{ employment\_questions.csv}

\CommentTok{\# Income questions (example)}
\FunctionTok{grep} \AttributeTok{{-}i} \StringTok{"income\textbackslash{}|earn\textbackslash{}|wage\textbackslash{}|salary"}\NormalTok{ expert\_review\_combined.csv }\OperatorTok{\textgreater{}}\NormalTok{ income\_questions.csv}
\end{Highlighting}
\end{Shaded}

\begin{center}\rule{0.5\linewidth}{0.5pt}\end{center}

\subsection{For Survey Designers}\label{for-survey-designers}

\textbf{Identifying Patterns}:

\begin{Shaded}
\begin{Highlighting}[]
\CommentTok{\# Group by barrier code to see common failure modes}
\FunctionTok{cut} \AttributeTok{{-}d}\StringTok{\textquotesingle{},\textquotesingle{}} \AttributeTok{{-}f9}\NormalTok{ expert\_review\_combined.csv }\KeywordTok{|} \FunctionTok{sort} \KeywordTok{|} \FunctionTok{uniq} \AttributeTok{{-}c} \KeywordTok{|} \FunctionTok{sort} \AttributeTok{{-}rn}

\CommentTok{\# Group by feasibility}
\FunctionTok{cut} \AttributeTok{{-}d}\StringTok{\textquotesingle{},\textquotesingle{}} \AttributeTok{{-}f7}\NormalTok{ expert\_review\_combined.csv }\KeywordTok{|} \FunctionTok{sort} \KeywordTok{|} \FunctionTok{uniq} \AttributeTok{{-}c}
\end{Highlighting}
\end{Shaded}

\textbf{Use Cases}: - \textbf{Questionnaire design}: Identify wording
patterns that enable/prevent harmonization - \textbf{New question
development}: Check if similar questions already harmonized -
\textbf{Cross-survey integration}: Plan data linkage based on
consolidable variables

\begin{center}\rule{0.5\linewidth}{0.5pt}\end{center}

\section{Summary Statistics}\label{summary-statistics-2}

\subsection{Overall Distribution}\label{overall-distribution}

\begin{longtable}[]{@{}llll@{}}
\toprule\noalign{}
Category & CPS & FoodAPS & Combined \\
\midrule\noalign{}
\endhead
\bottomrule\noalign{}
\endlastfoot
\textbf{Total} & 240 & 140 & 380 \\
\textbf{F1} & 37 (15.4\%) & 23 (16.4\%) & 60 (15.8\%) \\
\textbf{F2} & 63 (26.3\%) & 45 (32.1\%) & 108 (28.4\%) \\
\textbf{F3} & 140 (58.3\%) & 72 (51.4\%) & 212 (55.8\%) \\
\textbf{Consolidable} & 100 (41.7\%) & 68 (48.6\%) & 168 (44.2\%) \\
\end{longtable}

\subsection{Triage Distribution}\label{triage-distribution}

\begin{longtable}[]{@{}llll@{}}
\toprule\noalign{}
Quadrant & CPS & FoodAPS & Combined \\
\midrule\noalign{}
\endhead
\bottomrule\noalign{}
\endlastfoot
\textbf{Q1} (Auto-accept) & 94 (39.2\%) & 57 (40.7\%) & 151 (39.7\%) \\
\textbf{Q2} (Auto-reject) & 84 (35.0\%) & 52 (37.1\%) & 136 (35.8\%) \\
\textbf{Q3} (Priority review) & 26 (10.8\%) & 14 (10.0\%) & 40
(10.5\%) \\
\textbf{Q4} (Secondary review) & 36 (15.0\%) & 17 (12.1\%) & 53
(13.9\%) \\
\end{longtable}

\subsection{Expert Review Load}\label{expert-review-load-1}

\begin{longtable}[]{@{}lll@{}}
\toprule\noalign{}
Category & Questions & \% of Total \\
\midrule\noalign{}
\endhead
\bottomrule\noalign{}
\endlastfoot
\textbf{Auto-processed} (Q1 + Q2) & 287 & 75.5\% \\
\textbf{Expert review} (Q3 + Q4) & 93 & 24.5\% \\
\textbf{Priority review} (Q3) & 40 & 10.5\% \\
\end{longtable}

\begin{center}\rule{0.5\linewidth}{0.5pt}\end{center}

\section{Validation Checklist}\label{validation-checklist}

When validating results, check:

\begin{itemize}
\tightlist
\item[$\square$]
  All 380 questions have exactly one best match
\item[$\square$]
  F1/F2 questions have \texttt{has\_consolidable\_path\ =\ True}
\item[$\square$]
  F3 questions have valid barrier code (CC.\emph{, TC.}, RS.\emph{,
  PC.}, MC.\emph{, PM.})
\item[$\square$]
  Borda scores in {[}0, 1{]}
\item[$\square$]
  Entropy scores in {[}0, 1{]}
\item[$\square$]
  Triage quadrant matches Borda-Entropy rules
\item[$\square$]
  Reasoning summary provided for all questions
\item[$\square$]
  No missing data in required fields
\end{itemize}

\begin{center}\rule{0.5\linewidth}{0.5pt}\end{center}

\section{Contact Information}\label{contact-information}

For questions about expert review tables: - \textbf{Technical issues}:
See \texttt{docs/SOFTWARE.md} for script documentation -
\textbf{Methodology questions}: See this report's Methodology section -
\textbf{Data corrections}: Submit corrections via project repository

\begin{center}\rule{0.5\linewidth}{0.5pt}\end{center}

\section{Next Steps}\label{next-steps-1}

\begin{enumerate}
\def\labelenumi{\arabic{enumi}.}
\tightlist
\item
  \textbf{Expert validation}: Complete Q3/Q4 review
\item
  \textbf{Quality assessment}: Evaluate classification accuracy
\item
  \textbf{Feedback integration}: Refine classifications based on expert
  input
\item
  \textbf{Consolidation pilots}: Test F1 recommendations in practice
\item
  \textbf{Documentation updates}: Incorporate lessons learned
\end{enumerate}

\begin{center}\rule{0.5\linewidth}{0.5pt}\end{center}

\textbf{Files Ready for Review}: -
\texttt{output/analysis/expert\_review\_cps.csv} (240 rows) -
\texttt{output/analysis/expert\_review\_foodaps.csv} (140 rows) -
\texttt{output/analysis/expert\_review\_combined.csv} (380 rows)

\textbf{Status}: ✅ Ready for expert validation




\end{document}
